\chapter{Cotangent Complexes and Spectral Deformation Theory}
\label{chap:spectral-deformation-theory}

\section{Purpose and standing conventions}
\label{sec:spectral-deformation-purpose}

This chapter develops the infinitesimal geometry of the spectral schemes
constructed in \Ref{chap:spectral-algebraic-geometry}. All affine,
Zariski, quasi-coherent, relative-spectrum, open--closed, and linear
geometric constructions are those of that chapter. The new input is the
tangent theory of commutative ring spectra and its globalization to the
existing structure sheaves of spectral schemes.

The chapter proceeds from the stabilization of commutative algebra to
split and nonsplit square-zero extensions, relative cotangent complexes,
multiplicative Postnikov invariants, affine and global obstruction theory,
and quasi-smooth closed immersions. The affine foundations are taken from
\cite[Sections~7.3--7.4]{LurieHigherAlgebra} and
\cite[Sections~1--3]{LurieDAGIV}. The sheaf-level cotangent and
square-zero constructions are the specialization to spectral schemes of
the cotangent theory of spectrally ringed $\infty$-topoi developed in
\cite[Chapter~17]{LurieSAG}. The global deformation statements below are
obtained by identifying an infinitesimal deformation with a square-zero
deformation of the structure sheaf on the fixed Zariski $\infty$-topos
and applying the sheaf-level tangent correspondence; they agree with the
spectral Deligne--Mumford formulation of
\cite[Section~19.4.3]{LurieSAG}. Closed immersions are those of
\Ref{def:affine-closed-immersion-spectral} and
\cite{LurieDAGIX}. The terminology and local geometry of quasi-smooth
closed immersions are standard in derived algebraic geometry; see
\cite[Section~2.1.2]{KhanQuasiSmoothBlowups}.

All conventions of \Ref{chap:spectral-algebraic-geometry} remain in
force. In particular, coordinate algebras and spectral schemes are
connective, while module categories and mapping spectra are taken in their
full stable \(\infty\)-categories. For a commutative ring spectrum \(A\),
put
\[
    L_A
    :=
    L_{A/\mathbb S}.
\]
For a spectral scheme \(X\), put
\[
    L_X
    :=
    L_{X/\Spec(\mathbb S)}.
\]

A square-zero extension used in connective spectral geometry will always
carry its classifying derivation and its specified pullback
identification. This structured data is retained because an underlying
morphism of commutative ring spectra does not, in unrestricted
connectivity, determine the coefficient module or the classifying
derivation.

When no relative base is displayed, a structured square-zero extension
or structured infinitesimal extension is understood to be relative to
the sphere spectrum, respectively to
\[
    \Spec(\mathbb S).
\]

\section{Tangent theory and spectral derivations}
\label{sec:spectral-tangent-derivations}

\subsection{The tangent category of commutative algebra}

\begin{theorem}[Tangent category of a commutative ring spectrum]
\label{thm:tangent-category-calg}
Let \(B\in\CAlg(\Sp)\). There is a canonical equivalence of stable
\(\infty\)-categories
\[
    \operatorname{Sp}
    \bigl(
        \CAlg(\Sp)_{/B}
    \bigr)
    \simeq
    \Mod_B.
\]
The zeroth-space functor has target
\[
    \Omega^\infty
    \colon
    \operatorname{Sp}
    \bigl(
        \CAlg(\Sp)_{/B}
    \bigr)
    \longrightarrow
    \bigl(
        \CAlg(\Sp)_{/B}
    \bigr)_*,
\]
where
\[
    \bigl(
        \CAlg(\Sp)_{/B}
    \bigr)_*
    \simeq
    \CAlg(\Sp)_{B//B}
\]
is the \(\infty\)-category of diagrams
\[
    B
    \longrightarrow
    C
    \longrightarrow
    B
\]
whose composite is equivalent to \(\id_B\).

Under the tangent equivalence, the zeroth-space functor carries a
\(B\)-module \(M\) to the split square-zero algebra
\[
    B
    \longrightarrow
    B\oplus M
    \longrightarrow
    B.
\]
\end{theorem}

\begin{proof}
This is the commutative specialization of the tangent-category theorem
for algebras over a coherent \(\infty\)-operad; see
\cite[Theorem~7.3.4.13, Corollary~7.3.4.14, and
Remark~7.3.4.15]{LurieHigherAlgebra}.

The zeroth-space functor from the stabilization of an
\(\infty\)-category with finite limits lands in the pointed objects of
that \(\infty\)-category. A pointed object of
\(\CAlg(\Sp)_{/B}\) consists of an augmented commutative algebra
\[
    C
    \longrightarrow
    B
\]
together with a morphism from the terminal object \(B\to B\), hence a
section
\[
    B
    \longrightarrow
    C.
\]
It is therefore equivalently a diagram
\[
    B
    \longrightarrow
    C
    \longrightarrow
    B
\]
whose composite is equivalent to \(\id_B\).
\end{proof}

\subsection{Split square-zero algebras}

\begin{construction}[Split square-zero algebra]
\label{con:split-square-zero-algebra}
For a commutative ring spectrum \(B\), write
\[
    \CAlg(\Sp)_{B//B}
\]
for the \(\infty\)-category of diagrams
\[
    B
    \longrightarrow
    C
    \longrightarrow
    B
\]
whose composite is equivalent to \(\id_B\).

Let \(M\in\Mod_B\). The tangent equivalence of
\Ref{thm:tangent-category-calg}, followed by the zeroth-space functor,
supplies a commutative ring spectrum
\[
    B\oplus M
\]
together with morphisms
\[
    B
    \xrightarrow{s}
    B\oplus M
    \xrightarrow{p}
    B
\]
satisfying
\[
    p\circ s
    \simeq
    \id_B.
\]
Its underlying \(B\)-module is the direct sum \(B\oplus M\), and the
multiplication on the \(M\)-summand is coherently null. It is called the
\textbf{split square-zero algebra of \(B\) by \(M\)}.
\end{construction}

\begin{proposition}[Functoriality of split square-zero algebras]
\label{prop:split-square-zero-functoriality}
For every commutative ring spectrum \(B\), the construction
\[
    M
    \longmapsto
    \bigl(
        B
        \longrightarrow
        B\oplus M
        \longrightarrow
        B
    \bigr)
\]
defines a limit-preserving functor
\[
    \Mod_B
    \longrightarrow
    \CAlg(\Sp)_{B//B}.
\]
If \(B\) and \(M\) are connective, then \(B\oplus M\) is connective.
\end{proposition}

\begin{proof}
Under the equivalence
\[
    \Mod_B
    \simeq
    \operatorname{Sp}
    \bigl(
        \CAlg(\Sp)_{/B}
    \bigr),
\]
the displayed functor is the zeroth-space functor
\[
    \Omega^\infty
    \colon
    \operatorname{Sp}
    \bigl(
        \CAlg(\Sp)_{/B}
    \bigr)
    \longrightarrow
    \bigl(
        \CAlg(\Sp)_{/B}
    \bigr)_*
    \simeq
    \CAlg(\Sp)_{B//B}.
\]
The zeroth-space functor is a right adjoint and therefore preserves
limits.

The underlying spectrum of \(B\oplus M\) is the direct sum of the
underlying spectra of \(B\) and \(M\). Hence \(B\oplus M\) is connective
whenever \(B\) and \(M\) are connective. See
\cite[Remark~7.3.4.15]{LurieHigherAlgebra}.
\end{proof}

\subsection{Relative derivations}

\begin{definition}[Spectral derivation]
\label{def:spectral-derivation}
Let $A\to B$ be a morphism of commutative ring spectra and let
$M\in\Mod_B$. The \textbf{space of $A$-linear derivations from $B$ to
$M$} is
\[
    \operatorname{Der}_A(B,M)
    :=
    \Map_{\CAlg(\Sp)_{A//B}}
    \bigl(
        B,
        B\oplus M
    \bigr).
\]
Thus an $A$-linear derivation is a section of
\[
    B\oplus M
    \longrightarrow
    B
\]
in commutative $A$-algebras.
\end{definition}

\begin{proposition}[Elementary functoriality of derivations]
\label{prop:elementary-derivation-functoriality}
Let
\[
    A
    \longrightarrow
    B
    \longrightarrow
    C
\]
be morphisms of commutative ring spectra.
\begin{enumerate}[label=(\roman*)]
    \item A morphism of $B$-modules $M\to M'$ induces a natural morphism
    \[
        \operatorname{Der}_A(B,M)
        \longrightarrow
        \operatorname{Der}_A(B,M').
    \]

    \item For every $C$-module $N$, restriction along $B\to C$ induces a
    natural morphism
    \[
        \operatorname{Der}_A(C,N)
        \longrightarrow
        \operatorname{Der}_A(B,N),
    \]
    where $N$ on the right is regarded as a $B$-module.
\end{enumerate}
\end{proposition}

\begin{proof}
Part (i) follows by postcomposition with the induced morphism
$B\oplus M\to B\oplus M'$. For part (ii), compose an $A$-algebra
section $C\to C\oplus N$ with $B\to C$ and use the canonical pullback
identification
\[
    B\times_C(C\oplus N)
    \simeq
    B\oplus N.
\]
\end{proof}

\section{Cotangent complexes of commutative ring spectra}
\label{sec:affine-cotangent-complexes}

\subsection{Existence and the universal derivation}

\begin{theorem}[Existence of the relative cotangent complex]
\label{thm:relative-cotangent-existence}
Let $A\to B$ be a morphism of commutative ring spectra. There exists a
$B$-module
\[
    L_{B/A}
    \in
    \Mod_B
\]
and a universal $A$-linear derivation such that, for every $B$-module
$M$, there is a natural equivalence
\[
    \Map_{\Mod_B}
    \bigl(
        L_{B/A},
        M
    \bigr)
    \simeq
    \operatorname{Der}_A(B,M).
\]
The construction is functorial in commutative triangles of ring spectra.
\end{theorem}

\begin{proof}
The relative cotangent complex is the relative cotangent object in the
tangent bundle of $\CAlg(\Sp)$. The tangent equivalence of
\Ref{thm:tangent-category-calg} identifies its fibre at $B$ with
$\Mod_B$, and the stabilization adjunction gives the displayed
representing equivalence. See
\cite[Sections~7.3.2--7.3.4]{LurieHigherAlgebra} and
\cite[Section~1]{LurieDAGIV}.
\end{proof}

\begin{corollary}[Cotangent complex of an equivalence]
\label{cor:cotangent-equivalence-zero}
If $A\to B$ is an equivalence, then
\[
    L_{B/A}
    \simeq
    0.
\]
\end{corollary}

\begin{proof}
The space of relative derivations is contractible for every coefficient
module, so its representing object is zero.
\end{proof}

\subsection{Free commutative algebras}

\begin{proposition}[Cotangent complex of a free algebra]
\label{prop:cotangent-free-algebra}
Let $A\in\CAlg(\Sp)$ and let $M\in\Mod_A$. Put
\[
    B
    :=
    \operatorname{Sym}_A(M).
\]
Then there is a natural equivalence of $B$-modules
\[
    L_{B/A}
    \simeq
    B\otimes_A M.
\]
\end{proposition}

\begin{proof}
For every $B$-module $N$, the free-algebra and extension-of-scalars
adjunctions give
\[
\begin{aligned}
    \operatorname{Der}_A(B,N)
    &\simeq
    \Map_{\Mod_A}(M,N)
    \\
    &\simeq
    \Map_{\Mod_B}(B\otimes_A M,N).
\end{aligned}
\]
Yoneda and \Ref{thm:relative-cotangent-existence} give the result.
See also \cite[Proposition~7.4.3.14]{LurieHigherAlgebra}.
\end{proof}

\begin{corollary}[Cotangent complex of brave new affine space]
\label{cor:cotangent-brave-affine-space}
For
\[
    A\{t_1,\ldots,t_n\}
    :=
    \operatorname{Sym}_A(A^{\oplus n}),
\]
there is a natural equivalence
\[
    L_{A\{t_1,\ldots,t_n\}/A}
    \simeq
    A\{t_1,\ldots,t_n\}^{\oplus n}.
\]
\end{corollary}

\begin{proof}
Apply \Ref{prop:cotangent-free-algebra} to $M=A^{\oplus n}$.
\end{proof}

\subsection{Base change and transitivity}

\begin{theorem}[Base change for cotangent complexes]
\label{thm:cotangent-base-change}
Let
\[
\begin{tikzcd}[column sep=large, row sep=large]
    A
        \arrow[r]
        \arrow[d]
    &
    B
        \arrow[d]
    \\
    A'
        \arrow[r]
    &
    B':=B\otimes_AA'
\end{tikzcd}
\]
be a pushout square of commutative ring spectra. Then there is a natural
equivalence of $B'$-modules
\[
    L_{B/A}\otimes_BB'
    \simeq
    L_{B'/A'}.
\]
\end{theorem}

\begin{proof}
For every $B'$-module $N$, extension and restriction of scalars give
\[
\begin{aligned}
    \Map_{\Mod_{B'}}(L_{B/A}\otimes_BB',N)
    &\simeq
    \Map_{\Mod_B}(L_{B/A},N)
    \\
    &\simeq
    \operatorname{Der}_A(B,N)
    \\
    &\simeq
    \operatorname{Der}_{A'}(B',N)
    \\
    &\simeq
    \Map_{\Mod_{B'}}(L_{B'/A'},N).
\end{aligned}
\]
The middle equivalence is the pushout universal property. Yoneda gives
the result. See
\cite[Proposition~7.3.3.7]{LurieHigherAlgebra}.
\end{proof}

\begin{corollary}[Base change of derivations]
\label{cor:derivation-base-change}
In the situation of \Ref{thm:cotangent-base-change}, for every
$B'$-module $N$ there is a natural equivalence
\[
    \operatorname{Der}_{A'}(B',N)
    \simeq
    \operatorname{Der}_A(B,N).
\]
\end{corollary}

\begin{proof}
Apply the representing property and
\Ref{thm:cotangent-base-change}.
\end{proof}

\begin{theorem}[Transitivity triangle]
\label{thm:cotangent-transitivity}
For composable morphisms of commutative ring spectra
\[
    A
    \longrightarrow
    B
    \longrightarrow
    C,
\]
there is a natural cofibre sequence of $C$-modules
\[
    L_{B/A}\otimes_BC
    \longrightarrow
    L_{C/A}
    \longrightarrow
    L_{C/B}.
\]
\end{theorem}

\begin{proof}
Restriction of an $A$-linear derivation of $C$ to $B$ gives the
corresponding fibre sequence of derivation functors. Representing those
functors in the stable category $\Mod_C$ gives the asserted cofibre
sequence. See
\cite[Proposition~7.3.3.5 and Corollary~7.3.3.6]{LurieHigherAlgebra}.
\end{proof}

\subsection{Connectivity and classical comparison}

\begin{proposition}[Connectivity and degree-zero comparison]
\label{prop:cotangent-connectivity-degree-zero}
Let $A\to B$ be a morphism of connective commutative ring spectra. Then
\[
    L_{B/A}
    \in
    \Mod_{B,\geq0},
\]
and there is a natural isomorphism of $\pi_0(B)$-modules
\[
    \pi_0(L_{B/A})
    \simeq
    \Omega^1_{\pi_0(B)/\pi_0(A)}.
\]
\end{proposition}

\begin{proof}
This is
\cite[Proposition~7.4.3.9]{LurieHigherAlgebra}; see also
\cite[Section~2]{LurieDAGIV}.
\end{proof}

\begin{theorem}[Derived conormal comparison]
\label{thm:derived-conormal-comparison}
Let
\[
    f
    \colon
    A
    \longrightarrow
    B
\]
be a morphism of connective commutative ring spectra which is surjective
on \(\pi_0\). Put
\[
    K
    :=
    \operatorname{fib}(f).
\]
Then \(K\) is connective, \(L_{B/A}\) is \(1\)-connective, and there is a
canonical morphism of connective \(B\)-modules
\[
    B\otimes_AK
    \longrightarrow
    L_{B/A}[-1]
\]
which induces an isomorphism
\[
    \pi_0(B\otimes_AK)
    \xrightarrow{\simeq}
    \pi_0
    \bigl(
        L_{B/A}[-1]
    \bigr).
\]

The comparison is natural in commutative triangles. More precisely, let
\[
    A
    \longrightarrow
    C
    \longrightarrow
    B
\]
be a commutative triangle of connective commutative ring spectra such
that both
\[
    A
    \longrightarrow
    C
    \qquad\text{and}\qquad
    A
    \longrightarrow
    B
\]
are surjective on \(\pi_0\). Put
\[
    K_C
    :=
    \operatorname{fib}(A\to C),
    \qquad
    K_B
    :=
    \operatorname{fib}(A\to B).
\]
Then there is a natural commutative square of \(B\)-modules
\[
\begin{tikzcd}[column sep=large, row sep=large]
    B\otimes_AK_C
        \arrow[r]
        \arrow[d]
    &
    \bigl(
        L_{C/A}\otimes_CB
    \bigr)[-1]
        \arrow[d]
    \\
    B\otimes_AK_B
        \arrow[r]
    &
    L_{B/A}[-1].
\end{tikzcd}
\]
The horizontal morphisms are the derived-conormal comparison morphisms,
with the upper one extended from \(C\) to \(B\), and the right vertical
morphism is induced by transitivity.
\end{theorem}

\begin{proof}
Since \(A\to B\) is surjective on \(\pi_0\), its cofibre is
\(1\)-connective. In the stable category of spectra there is a canonical
equivalence
\[
    \operatorname{cofib}(A\to B)
    \simeq
    K[1].
\]
The canonical morphism
\[
    B\otimes_A
    \operatorname{cofib}(A\to B)
    \longrightarrow
    L_{B/A}
\]
has \(2\)-connective fibre by the relative cotangent connectivity
theorem. After using
\[
    \operatorname{cofib}(A\to B)
    \simeq
    K[1]
\]
and shifting by \([-1]\), one obtains a canonical morphism
\[
    B\otimes_AK
    \longrightarrow
    L_{B/A}[-1]
\]
which induces an isomorphism on \(\pi_0\). The same connectivity theorem
shows that \(L_{B/A}\) is \(1\)-connective. The fibre \(K\) is connective
because \(A\) and \(B\) are connective and \(A\to B\) is surjective on
\(\pi_0\).

The morphism from the relative cofibre to the relative cotangent complex
is natural in commutative diagrams of commutative ring spectra. Applying
this naturality to the triangle
\[
    A
    \longrightarrow
    C
    \longrightarrow
    B,
\]
then identifying the relevant cofibres with \(K_C[1]\) and \(K_B[1]\),
extending scalars from \(C\) to \(B\), and shifting by \([-1]\), gives
the displayed commutative square. The right vertical morphism agrees
with the morphism induced by the transitivity triangle.

See
\cite[Theorem~7.4.3.1 and Corollary~7.4.3.2]{LurieHigherAlgebra}.
\end{proof}

\begin{lemma}[Classical quotient of the derived conormal module]
\label{lem:derived-conormal-classical-quotient}
In the situation of \Ref{thm:derived-conormal-comparison}, put
\[
    R:=\pi_0(A),
    \qquad
    S:=\pi_0(B),
    \qquad
    I:=\ker(R\to S).
\]
Then the morphism $K\to A$ induces a natural surjection
\[
    \pi_0(B\otimes_AK)
    \longrightarrow
    I/I^2.
\]
\end{lemma}

\begin{proof}
Since $B$ and $K$ are connective, the connective tensor-product spectral
sequence gives
\[
    \pi_0(B\otimes_AK)
    \simeq
    S\otimes_R\pi_0(K).
\]
The long exact sequence of the fibre $K\to A\to B$ shows that
$\pi_0(K)\to R$ has image $I$ and is surjective onto $I$. Tensoring
with $S$ over $R$ gives a surjection
\[
    S\otimes_R\pi_0(K)
    \longrightarrow
    S\otimes_RI
    \simeq
    I/I^2.
\]
\end{proof}

\begin{remark}[Classical and derived conormal modules]
\label{rem:classical-derived-conormal-distinction}
Let
\[
    I
    :=
    \ker
    \bigl(
        \pi_0(A)
        \longrightarrow
        \pi_0(B)
    \bigr).
\]
The spectral conormal module
\[
    \pi_0\bigl(L_{B/A}[-1]\bigr)
\]
need not be isomorphic to $I/I^2$. The derived fibre $K$ can carry
conormal information which is invisible in the ordinary ideal $I$.
Accordingly, local equations for a spectral closed immersion must be
lifted from the derived fibre; such a lift includes a chosen
nullhomotopy after restriction to $B$.
\end{remark}

\begin{corollary}[Cotangent rigidity]
\label{cor:cotangent-rigidity}
Let $A\to B$ be a morphism of connective commutative ring spectra. If
\[
    \pi_0(A)
    \xrightarrow{\simeq}
    \pi_0(B)
\]
and
\[
    L_{B/A}
    \simeq
    0,
\]
then $A\to B$ is an equivalence.
\end{corollary}

\begin{proof}
This is
\cite[Corollary~7.4.3.4]{LurieHigherAlgebra}; see also
\cite[Corollary~2.7]{LurieDAGIV}.
\end{proof}

\subsection{Finiteness detected by the cotangent complex}

\begin{theorem}[Cotangent criterion for finite presentation]
\label{thm:cotangent-finite-presentation}
Let $A\to B$ be a morphism of connective commutative ring spectra.
\begin{enumerate}[label=(\roman*)]
    \item If $B$ is of finite presentation over $A$, then $L_{B/A}$ is a
    perfect $B$-module.

    \item Conversely, if $L_{B/A}$ is perfect and $\pi_0(B)$ is finitely
    presented as a $\pi_0(A)$-algebra, then $B$ is of finite presentation
    over $A$.
\end{enumerate}
\end{theorem}

\begin{proof}
This is
\cite[Theorem~7.4.3.18]{LurieHigherAlgebra}.
\end{proof}

\section{Structured square-zero extensions and Postnikov theory}
\label{sec:structured-square-zero-postnikov}

\subsection{The extension classified by a derivation}

\begin{definition}[Structured square-zero extension]
\label{def:structured-square-zero-extension}
Let \(A\to B\) be a morphism of commutative ring spectra and let
\(M\in\Mod_B\). A morphism
\[
    \eta
    \colon
    L_{B/A}
    \longrightarrow
    M[1]
\]
represents an \(A\)-linear derivation
\[
    d_\eta
    \colon
    B
    \longrightarrow
    B\oplus M[1].
\]
Let
\[
    d_0
    \colon
    B
    \longrightarrow
    B\oplus M[1]
\]
denote the zero derivation. Define
\[
    B^\eta
    :=
    B
    \times_{B\oplus M[1]}
    B,
\]
where the right vertical morphism in the pullback square below is
\(d_\eta\) and the lower horizontal morphism is \(d_0\):
\[
\begin{tikzcd}[column sep=large, row sep=large]
    B^\eta
        \arrow[r, "q_\eta"]
        \arrow[d, "r_\eta"']
    &
    B
        \arrow[d, "d_\eta"]
    \\
    B
        \arrow[r, "d_0"']
    &
    B\oplus M[1].
\end{tikzcd}
\]
The derivation \(\eta\), the object \(B^\eta\), and the specified
pullback identification form the
\textbf{structured square-zero extension of \(B\) by \(M\) over \(A\)
classified by \(\eta\)}.

By convention, the morphism
\[
    q_\eta
    \colon
    B^\eta
    \longrightarrow
    B
\]
is the underlying square-zero extension. Thus every subsequent reference
to the morphism \(B^\eta\to B\) means \(q_\eta\).
\end{definition}

\begin{proposition}[Fibre of a structured square-zero extension]
\label{prop:fibre-structured-square-zero}
In the situation of \Ref{def:structured-square-zero-extension}, the
fibre of
\[
    q_\eta
    \colon
    B^\eta
    \longrightarrow
    B
\]
is canonically equivalent to \(M\), regarded as a
\(B^\eta\)-module by restriction of scalars along \(q_\eta\). The induced
multiplication on the fibre is null.

If \(B\) and \(M\) are connective, then \(B^\eta\) is connective and
\[
    \pi_0(q_\eta)
    \colon
    \pi_0(B^\eta)
    \longrightarrow
    \pi_0(B)
\]
is surjective with square-zero kernel.
\end{proposition}
\begin{proof}
The forgetful functor
\[
    \CAlg(\Sp)
    \longrightarrow
    \Sp
\]
creates limits. Hence the underlying square of spectra
\[
\begin{tikzcd}[column sep=large, row sep=large]
    B^\eta
        \arrow[r, "q_\eta"]
        \arrow[d, "r_\eta"']
    &
    B
        \arrow[d, "d_\eta"]
    \\
    B
        \arrow[r, "d_0"']
    &
    B\oplus M[1]
\end{tikzcd}
\]
is a pullback square.

As a morphism of underlying spectra, the zero derivation
\[
    d_0
    \colon
    B
    \longrightarrow
    B\oplus M[1]
\]
is the inclusion of the first summand. Therefore
\[
    \operatorname{cofib}(d_0)
    \simeq
    M[1].
\]
By stability,
\[
    \operatorname{fib}(d_0)
    \simeq
    M.
\]
Fibres are preserved by pullback, so
\[
    \operatorname{fib}(q_\eta)
    \simeq
    M.
\]
The resulting equivalence is an equivalence of
\(B^\eta\)-modules, where \(M\) is regarded as a
\(B^\eta\)-module by restriction of scalars along
\[
    q_\eta
    \colon
    B^\eta
    \longrightarrow
    B.
\]
Thus there is a canonical fibre sequence of
\(B^\eta\)-modules
\[
    M
    \longrightarrow
    B^\eta
    \xrightarrow{q_\eta}
    B.
\]

The fibre of a morphism of commutative ring spectra carries a canonical
nonunital commutative algebra structure over the source. In the present
case, \(q_\eta\) is the extension obtained from the derivation
\(d_\eta\) by the tangent-theoretic square-zero construction. Under the
tangent equivalence of
\Ref{thm:tangent-category-calg}, its coefficient nonunital algebra is
the \(B^\eta\)-module \(M\) equipped with the zero multiplication.
Equivalently, the binary operation
\[
    M\otimes_{B^\eta}M
    \longrightarrow
    M
\]
is coherently null. This is the square-zero property of the
derivation-to-extension construction.

Suppose now that \(B\) and \(M\) are connective. Since the connective
part of \(\Mod_{B^\eta}\) is closed under extensions, the fibre sequence
\[
    M
    \longrightarrow
    B^\eta
    \longrightarrow
    B
\]
shows that \(B^\eta\) is connective.

The associated long exact sequence of homotopy groups contains
\[
    \pi_0(B^\eta)
    \xrightarrow{\pi_0(q_\eta)}
    \pi_0(B)
    \longrightarrow
    \pi_{-1}(M).
\]
Since \(M\) is connective,
\[
    \pi_{-1}(M)
    =
    0.
\]
Therefore
\[
    \pi_0(q_\eta)
    \colon
    \pi_0(B^\eta)
    \longrightarrow
    \pi_0(B)
\]
is surjective.

Its kernel is the image of
\[
    \pi_0(M)
    \longrightarrow
    \pi_0(B^\eta).
\]
Because the multiplication
\[
    M\otimes_{B^\eta}M
    \longrightarrow
    M
\]
is coherently null, the product of any two elements in this image is
zero. Hence the kernel of \(\pi_0(q_\eta)\) is a square-zero ideal.
\end{proof}

\subsection{Classification and recognition}

\begin{theorem}[Structured square-zero classification]
\label{thm:structured-square-zero-classification}
Let $A\to B$ be a morphism of commutative ring spectra and let
$M\in\Mod_B$.

Let
\[
    \operatorname{SqZ}^{\mathrm{str}}_A(B;M)
\]
denote the $\infty$-groupoid whose objects are structured square-zero
extensions of $B$ by $M$ over $A$, consisting of:
\begin{enumerate}[label=(\roman*)]
    \item a morphism
    \[
        \eta
        \colon
        L_{B/A}
        \longrightarrow
        M[1];
    \]

    \item a commutative $A$-algebra $B^\eta$ equipped with a morphism
    \[
        q_\eta
        \colon
        B^\eta
        \longrightarrow
        B;
    \]

    \item a specified pullback identification
    \[
        B^\eta
        \simeq
        B
        \times_{B\oplus M[1]}
        B
    \]
    formed using the derivation represented by $\eta$ and the zero
    derivation.
\end{enumerate}
Its morphisms are equivalences over $B$ which preserve the identified
coefficient module $M$, the classifying morphism $\eta$, and the
specified pullback identification, together with the required coherent
homotopies.

There are natural equivalences
\[
    \operatorname{SqZ}^{\mathrm{str}}_A(B;M)
    \simeq
    \operatorname{Der}_A(B,M[1])
    \simeq
    \Map_{\Mod_B}
    \bigl(
        L_{B/A},
        M[1]
    \bigr).
\]
\end{theorem}

\begin{proof}
The second equivalence is the universal property of the relative
cotangent complex:
\[
    \operatorname{Der}_A(B,M[1])
    \simeq
    \Map_{\Mod_B}
    \bigl(
        L_{B/A},
        M[1]
    \bigr).
\]

For a derivation
\[
    d_\eta
    \colon
    B
    \longrightarrow
    B\oplus M[1],
\]
the construction of
\Ref{def:structured-square-zero-extension} forms the pullback
\[
    B^\eta
    :=
    B
    \times_{B\oplus M[1]}
    B,
\]
where the second map from $B$ to $B\oplus M[1]$ is the zero
derivation. This produces a structured square-zero extension of $B$ by
$M$ over $A$.

Conversely, the retained classifying morphism and specified pullback
identification recover the corresponding derivation. For a fixed
derivation, the space of choices of a pullback object together with a
specified pullback identification is contractible. Equivalences between
such structured extensions are therefore exactly equivalences preserving
the classifying derivation, the coefficient identification, and the
pullback data.

Hence the square-zero construction induces the first displayed
equivalence. See
\cite[Section~7.4.1]{LurieHigherAlgebra} and
\cite[Section~3.1]{LurieDAGIV}.
\end{proof}

\begin{remark}[Recognition of underlying extensions]
\label{rem:recognition-underlying-square-zero}
The theorem classifies structured extensions. Recovering the coefficient
module and classifying derivation from an underlying morphism requires a
separate recognition theorem.

Lurie's \(n\)-small extension theorem supplies such a recognition for
the bounded extensions occurring in Postnikov towers: on the specified
connectivity and truncation range, the structured derivation is recovered
from the underlying extension. See
\cite[Theorem~3.17 and Corollaries~3.18--3.19]{LurieDAGIV} and
\cite[Section~7.4.1]{LurieHigherAlgebra}.

For the arbitrary connective square-zero extensions used in the
deformation theory below, the coefficient module, the classifying
derivation, and the specified pullback identification are retained as
part of the input. No reconstruction of this structured data from the
underlying morphism alone is asserted.
\end{remark}

\subsection{Multiplicative Postnikov extensions}

\begin{theorem}[Postnikov stages are square-zero extensions]
\label{thm:postnikov-square-zero}
Let $B$ be a connective commutative ring spectrum. For every integer
$n\geq1$, the truncation morphism
\[
    \tau_{\leq n}B
    \longrightarrow
    \tau_{\leq n-1}B
\]
is canonically a structured square-zero extension with fibre
\[
    H\pi_n(B)[n].
\]
The coefficient module is regarded as a module over
$\tau_{\leq n-1}B$ through
\[
    \tau_{\leq n-1}B
    \longrightarrow
    H\pi_0(B),
\]
and the extension is classified by a multiplicative Postnikov
$k$-invariant
\[
    L_{\tau_{\leq n-1}B}
    \longrightarrow
    H\pi_n(B)[n+1].
\]
\end{theorem}

\begin{proof}
The underlying truncation map has the displayed fibre. It is an
$n$-small extension, so the connective recognition theorem identifies it
with the square-zero extension classified by its multiplicative
$k$-invariant. See
\cite[Corollary~3.19]{LurieDAGIV} and
\cite[Corollary~7.4.1.28]{LurieHigherAlgebra}.
\end{proof}

\begin{corollary}[Relative Postnikov extensions]
\label{cor:relative-postnikov-square-zero}
Let $A\to B$ be a morphism of connective commutative ring spectra. For
every $n\geq1$, the map
\[
    \tau_{\leq n}B
    \longrightarrow
    \tau_{\leq n-1}B
\]
is canonically a structured square-zero extension in
$\CAlg(\Sp)^{\mathrm{cn}}_{A/}$ by $H\pi_n(B)[n]$. Its relative
classifying morphism is
\[
    L_{\tau_{\leq n-1}B/A}
    \longrightarrow
    H\pi_n(B)[n+1].
\]
\end{corollary}

\begin{proof}
The canonical structural morphism
\[
    A
    \longrightarrow
    \tau_{\leq n}B
\]
is a lift of
\[
    A
    \longrightarrow
    \tau_{\leq n-1}B
\]
through the Postnikov square-zero extension
\[
    \tau_{\leq n}B
    \longrightarrow
    \tau_{\leq n-1}B.
\]
Consequently, the absolute classifying morphism
\[
    L_{\tau_{\leq n-1}B}
    \longrightarrow
    H\pi_n(B)[n+1]
\]
comes equipped with a canonical nullhomotopy after precomposition with
\[
    L_A\otimes_A\tau_{\leq n-1}B
    \longrightarrow
    L_{\tau_{\leq n-1}B}.
\]
By the transitivity cofibre sequence
\[
    L_A\otimes_A\tau_{\leq n-1}B
    \longrightarrow
    L_{\tau_{\leq n-1}B}
    \longrightarrow
    L_{\tau_{\leq n-1}B/A},
\]
this structured nullhomotopy determines the relative classifying
morphism
\[
    L_{\tau_{\leq n-1}B/A}
    \longrightarrow
    H\pi_n(B)[n+1].
\]
The relative structured square-zero classification now gives the
assertion.
\end{proof}

\begin{corollary}[Infinitesimal Postnikov reconstruction]
\label{cor:postnikov-infinitesimal-reconstruction}
Every connective commutative ring spectrum is reconstructed from its
ordinary ring of components by an inverse tower of structured square-zero
extensions:
\[
    B
    \simeq
    \varprojlim_n
    \tau_{\leq n}B,
    \qquad
    \tau_{\leq0}B
    \simeq
    H\pi_0(B).
\]
\end{corollary}

\begin{proof}
Combine \Ref{prop:operadic-postnikov-truncations} with
\Ref{thm:postnikov-square-zero}.
\end{proof}

\subsection{Formal lifting across a fixed extension}

\begin{proposition}[Formal lifting across a square-zero extension]
\label{prop:formal-lifting-square-zero}
Let
\[
    A
    \longrightarrow
    B,
    \qquad
    A
    \longrightarrow
    C
\]
be morphisms of commutative ring spectra. Let
\[
    C'
    \longrightarrow
    C
\]
be a structured square-zero extension in commutative $A$-algebras by a
$C$-module $J$, classified by
\[
    \delta
    \colon
    L_{C/A}
    \longrightarrow
    J[1].
\]
For an $A$-algebra morphism
\[
    \varphi
    \colon
    B
    \longrightarrow
    C,
\]
define
\[
    o(\varphi)
    \colon
    L_{B/A}\otimes_BC
    \longrightarrow
    L_{C/A}
    \xrightarrow{\delta}
    J[1].
\]
Then the space of lifts of $\varphi$ to an $A$-algebra morphism
$B\to C'$ is naturally equivalent to the space of nullhomotopies of
$o(\varphi)$. Consequently:
\begin{enumerate}[label=(\roman*)]
    \item a lift exists if and only if
    \[
        [o(\varphi)]
        =
        0
    \]
    in
    \[
        \pi_0
        \Map_{\Mod_C}
        \bigl(
            L_{B/A}\otimes_BC,
            J[1]
        \bigr);
    \]

    \item when nonempty, the lift space is a torsor under
    \[
        \Omega^\infty
        \map_{\Mod_C}
        \bigl(
            L_{B/A}\otimes_BC,
            J
        \bigr);
    \]

    \item if $L_{B/A}\simeq0$, the lift space is contractible.
\end{enumerate}
\end{proposition}

\begin{proof}
Write the structured extension as the pullback
\[
\begin{tikzcd}[column sep=large, row sep=large]
    C'
        \arrow[r]
        \arrow[d]
    &
    C
        \arrow[d, "d_\delta"]
    \\
    C
        \arrow[r, "d_0"']
    &
    C\oplus J[1].
\end{tikzcd}
\]
Applying $\Map_{\CAlg(\Sp)_{A/}}(B,-)$ gives a pullback square of
mapping spaces. Taking the fibre over $\varphi$ identifies the lift space
with the space of homotopies between $d_\delta\circ\varphi$ and the zero
derivation. Naturality of cotangent complexes identifies that derivation
with $o(\varphi)$. The remaining assertions follow by taking the
connected components and loop action of the mapping spectrum.
\end{proof}

\section{Affine deformation and obstruction theory}
\label{sec:affine-spectral-deformation-theory}

\subsection{Deformations over an infinitesimal base}

\begin{definition}[Affine deformation over a square-zero extension]
\label{def:affine-deformation-square-zero}
Let
\[
    A'
    \longrightarrow
    A
\]
be a structured square-zero extension of connective commutative ring
spectra by a connective \(A\)-module \(I\). Let \(B\) be a connective
commutative \(A\)-algebra.

A \textbf{deformation of \(B\) over \(A'\)} is a pair
\[
    (B',\alpha)
\]
consisting of a connective commutative \(A'\)-algebra \(B'\) and an
equivalence
\[
    \alpha
    \colon
    B'\otimes_{A'}A
    \xrightarrow{\simeq}
    B
\]
of commutative \(A\)-algebras.

A morphism
\[
    (B'_1,\alpha_1)
    \longrightarrow
    (B'_2,\alpha_2)
\]
is an equivalence of commutative \(A'\)-algebras
\[
    u
    \colon
    B'_1
    \xrightarrow{\simeq}
    B'_2
\]
together with the specified coherent compatibility
\[
    \alpha_2
    \circ
    (u\otimes_{A'}A)
    \simeq
    \alpha_1.
\]
The resulting \(\infty\)-groupoid is denoted
\[
    \operatorname{Def}_{A'/A}(B).
\]
\end{definition}

\subsection{The deformation space as a homotopy fibre}

\begin{theorem}[Affine deformation fibre sequence]
\label{thm:affine-deformation-homotopy-fibre}
Let
\[
    A'
    \longrightarrow
    A
\]
be classified by
\[
    \eta_A
    \colon
    L_A
    \longrightarrow
    I[1],
\]
where $I$ is connective. Let $B$ be a connective commutative $A$-algebra
and put
\[
    J
    :=
    B\otimes_AI.
\]
Base change gives
\[
    \eta_B
    \colon
    L_A\otimes_AB
    \longrightarrow
    J[1].
\]
Then there is a natural equivalence
\[
    \operatorname{Def}_{A'/A}(B)
    \simeq
    \operatorname{fib}_{\eta_B}
    \left(
        \Map_{\Mod_B}(L_B,J[1])
        \longrightarrow
        \Map_{\Mod_B}(L_A\otimes_AB,J[1])
    \right).
\]
Equivalently, a deformation is a morphism
\[
    \widetilde\eta
    \colon
    L_B
    \longrightarrow
    J[1]
\]
together with a specified homotopy
\[
    \widetilde\eta|_{L_A\otimes_AB}
    \simeq
    \eta_B.
\]
\end{theorem}

\begin{proof}
The structured extension $A'\to A$ is classified by $\eta_A$.
Lurie's tangent correspondence identifies connective commutative
$A'$-algebras reducing to $B$ with derivations of $B$ compatible with
$\eta_A$, and identifies the associated square of commutative ring
spectra with a pushout square. Under the cotangent representing property,
compatibility is precisely the displayed homotopy fibre. See
\cite[Proposition~3.22 and Theorem~3.25]{LurieDAGIV} and
\cite[Section~7.4.2]{LurieHigherAlgebra}.
\end{proof}

\subsection{Obstruction, deformations, and automorphisms}

\begin{theorem}[Affine cotangent obstruction theory]
\label{thm:affine-cotangent-obstruction}
In the situation of
\Ref{thm:affine-deformation-homotopy-fibre}, the transitivity triangle
\[
    L_A\otimes_AB
    \longrightarrow
    L_B
    \longrightarrow
    L_{B/A}
\]
determines a canonical obstruction class
\[
    \operatorname{ob}(B/A')
    \in
    \pi_0
    \map_{\Mod_B}
    \bigl(
        L_{B/A},
        J[2]
    \bigr).
\]
It has the following properties.
\begin{enumerate}[label=(\roman*)]
    \item The obstruction vanishes if and only if $B$ admits a
    deformation over $A'$.

    \item When it vanishes, the deformation $\infty$-groupoid is a
    torsor under
    \[
        \Omega^\infty
        \map_{\Mod_B}
        \bigl(
            L_{B/A},
            J[1]
        \bigr).
    \]

    \item Let
    \[
        \xi
        \in
        \operatorname{Def}_{A'/A}(B)
    \]
    be a fixed deformation. Write
    \[
        \operatorname{Aut}(\xi)
        :=
        \Omega_\xi
        \operatorname{Def}_{A'/A}(B)
    \]
    for its automorphism space in the deformation $\infty$-groupoid.
    Thus its elements are automorphisms over $A'$ which preserve the
    specified identification of the reduction with $B$, together with
    the required coherent homotopy. There is a natural equivalence
    \[
        \operatorname{Aut}(\xi)
        \simeq
        \Omega^\infty
        \map_{\Mod_B}
        \bigl(
            L_{B/A},
            J
        \bigr).
    \]
\end{enumerate}
\end{theorem}

\begin{proof}
Applying the exact contravariant functor
\[
    \map_{\Mod_B}(-,J[1])
\]
to the transitivity triangle gives a fibre sequence of spectra
\[
\begin{aligned}
    \map_{\Mod_B}(L_{B/A},J[1])
    &\longrightarrow
    \map_{\Mod_B}(L_B,J[1])
    \\
    &\longrightarrow
    \map_{\Mod_B}(L_A\otimes_AB,J[1])
    \\
    &\xrightarrow{\partial}
    \map_{\Mod_B}(L_{B/A},J[2]).
\end{aligned}
\]
Define
\[
    \operatorname{ob}(B/A')
    :=
    \partial(\eta_B).
\]

By
\Ref{thm:affine-deformation-homotopy-fibre},
the deformation $\infty$-groupoid is the homotopy fibre over $\eta_B$
of
\[
    \Map_{\Mod_B}(L_B,J[1])
    \longrightarrow
    \Map_{\Mod_B}(L_A\otimes_AB,J[1]).
\]
The point $\eta_B$ lifts through this map if and only if its boundary
under $\partial$ is null. This proves \textup{(i)}.

When the homotopy fibre is nonempty, translation by a chosen point
identifies it as a torsor under the homotopy fibre over zero. By the
displayed fibre sequence, that homotopy fibre is
\[
    \Omega^\infty
    \map_{\Mod_B}
    \bigl(
        L_{B/A},
        J[1]
    \bigr).
\]
This proves \textup{(ii)}.

For a fixed deformation $\xi$, its automorphism space in the deformation
$\infty$-groupoid is the based loop space
\[
    \Omega_\xi
    \operatorname{Def}_{A'/A}(B).
\]
Using the torsor description and the canonical equivalence between the
loop spaces at any two points of a grouplike infinite loop space gives
\[
\begin{aligned}
    \operatorname{Aut}(\xi)
    &\simeq
    \Omega
    \Omega^\infty
    \map_{\Mod_B}
    \bigl(
        L_{B/A},
        J[1]
    \bigr)
    \\
    &\simeq
    \Omega^\infty
    \map_{\Mod_B}
    \bigl(
        L_{B/A},
        J
    \bigr).
\end{aligned}
\]
This proves \textup{(iii)}. See
\cite[Remark~3.23]{LurieDAGIV}.
\end{proof}

\section{Global cotangent complexes on spectral schemes}
\label{sec:global-cotangent-spectral-schemes}

\subsection{The Zariski ringed $\infty$-topos}

Every spectral scheme $X$ constructed in
\Ref{chap:spectral-algebraic-geometry} has a canonical Zariski
$\infty$-topos, denoted
\[
    X_{\mathrm{Zar}},
\]
and its existing structure sheaf determines a sheaf of connective
commutative ring spectra
\[
    \mathcal O_X
    \in
    \CAlg
    \bigl(
        \operatorname{Shv}_{\Sp}(X_{\mathrm{Zar}})
    \bigr).
\]
The category $\QCoh(X)$ of
\Ref{def:qcoh-spectral-prestack} identifies with the full subcategory of
quasi-coherent $\mathcal O_X$-modules in the sheaf-theoretic module
category. Thus the ringed-$\infty$-topos formalism used below is intrinsic
machinery attached to the spectral schemes already constructed; it does
not alter the affine or Zariski foundations of the preceding chapter.

\subsection{Definition and affine comparison}

\begin{definition}[Global relative cotangent complex]
\label{def:global-relative-cotangent}
Let
\[
    f
    \colon
    X
    \longrightarrow
    Y
\]
be a morphism of spectral schemes. The induced geometric morphism of
Zariski $\infty$-topoi carries a morphism of sheaves of commutative ring
spectra
\[
    f^{-1}\mathcal O_Y
    \longrightarrow
    \mathcal O_X.
\]
The \textbf{relative cotangent complex} of $f$ is the sheaf-level
cotangent complex
\[
    L_{X/Y}
    :=
    L_{\mathcal O_X/f^{-1}\mathcal O_Y}.
\]
\end{definition}

\begin{theorem}[Quasi-coherence and affine compatibility]
\label{thm:global-cotangent-affine-compatibility}
Let $f:X\to Y$ be a morphism of spectral schemes.
\begin{enumerate}[label=(\roman*)]
    \item The sheaf-level cotangent complex is quasi-coherent and hence
    determines an object
    \[
        L_{X/Y}
        \in
        \QCoh(X).
    \]

    \item Let
    \[
        V=\Spec(B)
        \longrightarrow
        X
    \]
    be an affine open and let
    \[
        U=\Spec(A)
        \longrightarrow
        Y
    \]
    be an affine open through which $V\to Y$ factors. Then there is a
    canonical equivalence
    \[
        L_{X/Y}|_V
        \simeq
        L_{B/A}.
    \]

    \item The global cotangent complex is objectwise connective:
    \[
        L_{X/Y}
        \in
        \QCoh(X)_{\geq0}.
    \]
\end{enumerate}
\end{theorem}
\begin{proof}
Form the sheaf-level relative cotangent complex
\[
    L_{\mathcal O_X/f^{-1}\mathcal O_Y}
\]
in the stable category of \(\mathcal O_X\)-modules on
\(X_{\mathrm{Zar}}\). The existence of this object, its compatibility
with inverse-image functors, and its relative derivation universal
property are the sheaf-level cotangent constructions of
\cite[Sections~17.1.1--17.1.2]{LurieSAG}.

Let
\[
    V=\Spec(B)
    \longrightarrow
    X
\]
be an affine open, and suppose that the composite
\[
    V
    \longrightarrow
    X
    \longrightarrow
    Y
\]
factors through an affine open
\[
    U=\Spec(A)
    \longrightarrow
    Y.
\]
Restricting the morphism of structure sheaves
\[
    f^{-1}\mathcal O_Y
    \longrightarrow
    \mathcal O_X
\]
to \(V_{\mathrm{Zar}}\) identifies it with the morphism of Zariski
structure sheaves associated to the commutative ring-spectrum morphism
\[
    A
    \longrightarrow
    B.
\]

Compatibility of sheaf-level cotangent complexes with inverse image
along the open immersion \(V\hookrightarrow X\) therefore gives
\[
    L_{\mathcal O_X/f^{-1}\mathcal O_Y}|_V
    \simeq
    L_{\mathcal O_V/\mathcal O_U|_V}.
\]
The sheaf-level derivation universal property on the affine spectral
scheme \(V=\Spec(B)\) identifies the right side with the quasi-coherent
sheaf associated to the affine cotangent complex
\[
    L_{B/A}.
\]
Thus there is a canonical equivalence
\[
    L_{\mathcal O_X/f^{-1}\mathcal O_Y}|_V
    \simeq
    \widetilde{L_{B/A}}.
\]

These affine identifications are compatible with further affine
restriction. Indeed, if
\[
    W=\Spec(C)
    \longrightarrow
    V
\]
is an affine open, then the comparison on \(W\) is obtained from that on
\(V\) by cotangent base change. Hence the affine cotangent complexes on
the members of an affine atlas and on all their iterated intersections
form a compatible object of the descent limit computing
\(\QCoh(X)\).

By the atlas equivalence of
\Ref{thm:qcoh-atlas-computation}, the sheaf-level cotangent complex is
therefore quasi-coherent and determines an object
\[
    L_{X/Y}
    \in
    \QCoh(X).
\]
This proves \textup{(i)}, and the affine calculation proves
\textup{(ii)}.

Finally, for every such affine \(V=\Spec(B)\), the morphism
\(A\to B\) is a morphism of connective commutative ring spectra.
Therefore
\[
    L_{B/A}
    \in
    \Mod_{B,\geq0}
\]
by \Ref{prop:cotangent-connectivity-degree-zero}. Objectwise
connectivity in \(\QCoh(X)\) is detected on an affine atlas, so
\[
    L_{X/Y}
    \in
    \QCoh(X)_{\geq0}.
\]
This proves \textup{(iii)}.
\end{proof}

\begin{remark}[Atlas computation]
\label{rem:global-cotangent-atlas-computation}
Under the atlas equivalence of \Ref{thm:qcoh-atlas-computation}, the
object $L_{X/Y}$ is represented by the compatible family of affine
cotangent complexes on the members of the atlas and all iterated
intersections. Thus affine descent computes the intrinsic global object,
but is not needed to construct it.
\end{remark}

\subsection{Global base change and transitivity}

\begin{theorem}[Global base change]
\label{thm:global-cotangent-base-change}
Let
\[
\begin{tikzcd}[column sep=large, row sep=large]
    X'
        \arrow[r, "g'"]
        \arrow[d, "f'"']
    &
    X
        \arrow[d, "f"]
    \\
    Y'
        \arrow[r, "g"']
    &
    Y
\end{tikzcd}
\]
be a Cartesian square of spectral schemes. Then there is a natural
equivalence
\[
    L_{X'/Y'}
    \simeq
    (g')^*L_{X/Y}.
\]
\end{theorem}

\begin{proof}
Sheaf-level cotangent complexes commute with inverse image and algebraic
base change; see \cite[Section~17.1.1]{LurieSAG}. The assertion can also
be checked on compatible affine opens using
\Ref{thm:cotangent-base-change} and
\Ref{thm:global-cotangent-affine-compatibility}.
\end{proof}

\begin{theorem}[Global transitivity]
\label{thm:global-cotangent-transitivity}
For composable morphisms
\[
    X
    \xrightarrow{f}
    Y
    \xrightarrow{g}
    Z,
\]
there is a natural cofibre sequence in $\QCoh(X)$
\[
    f^*L_{Y/Z}
    \longrightarrow
    L_{X/Z}
    \longrightarrow
    L_{X/Y}.
\]
\end{theorem}

\begin{proof}
This is the transitivity sequence for the morphisms of structure sheaves
on $X_{\mathrm{Zar}}$; see \cite[Section~17.1.1]{LurieSAG}. Affine
restriction recovers \Ref{thm:cotangent-transitivity}.
\end{proof}

\begin{corollary}[Cotangent complex of an open immersion]
\label{cor:cotangent-open-immersion}
Let
\[
    j
    \colon
    U
    \hookrightarrow
    X
\]
be an open immersion of spectral schemes. Then
\[
    L_{U/X}
    \simeq
    0.
\]
\end{corollary}

\begin{proof}
The structure sheaf of $U$ is the inverse image of the structure sheaf of
$X$. Hence the relative morphism of sheaves of commutative ring spectra is
an equivalence, and its cotangent complex vanishes. See
\cite[Section~17.1.1]{LurieSAG}.
\end{proof}

\section{Global structured square-zero extensions}
\label{sec:global-structured-square-zero}

\subsection{The sheaf-level extension}

\begin{construction}[Global square-zero extension]
\label{con:global-square-zero-extension}
Let
\[
    f
    \colon
    X
    \longrightarrow
    Y
\]
be a morphism of spectral schemes, let
\[
    \mathcal M
    \in
    \QCoh(X),
\]
and let
\[
    \eta
    \colon
    L_{X/Y}
    \longrightarrow
    \mathcal M[1].
\]

The structural morphism
\[
    f^{-1}\mathcal O_Y
    \longrightarrow
    \mathcal O_X
\]
makes both
\[
    \mathcal O_X
    \qquad\text{and}\qquad
    \mathcal O_X\oplus\mathcal M[1]
\]
commutative \(f^{-1}\mathcal O_Y\)-algebras. The augmentation
\[
    \mathcal O_X\oplus\mathcal M[1]
    \longrightarrow
    \mathcal O_X
\]
is the projection onto the first summand.

By the relative cotangent universal property, the morphism \(\eta\)
represents a derivation
\[
    d_\eta
    \colon
    \mathcal O_X
    \longrightarrow
    \mathcal O_X\oplus\mathcal M[1]
\]
in the \(\infty\)-category
\[
    \CAlg
    \bigl(
        \operatorname{Shv}_{\Sp}(X_{\mathrm{Zar}})
    \bigr)_{f^{-1}\mathcal O_Y//\mathcal O_X}.
\]
Thus \(d_\eta\) is \(f^{-1}\mathcal O_Y\)-linear. Equivalently, its
underlying absolute derivation is classified by the composite
\[
    L_X
    \longrightarrow
    L_{X/Y}
    \xrightarrow{\eta}
    \mathcal M[1].
\]

Let
\[
    d_0
    \colon
    \mathcal O_X
    \longrightarrow
    \mathcal O_X\oplus\mathcal M[1]
\]
denote the zero \(f^{-1}\mathcal O_Y\)-linear derivation.

Define
\[
    \mathcal O_X^\eta
    :=
    \mathcal O_X
    \times_{\mathcal O_X\oplus\mathcal M[1]}
    \mathcal O_X
\]
by the specified pullback square
\[
\begin{tikzcd}[column sep=large, row sep=large]
    \mathcal O_X^\eta
        \arrow[r, "q_\eta^\sharp"]
        \arrow[d, "r_\eta^\sharp"']
    &
    \mathcal O_X
        \arrow[d, "d_\eta"]
    \\
    \mathcal O_X
        \arrow[r, "d_0"']
    &
    \mathcal O_X\oplus\mathcal M[1].
\end{tikzcd}
\]
The pullback is formed in
\[
    \CAlg
    \bigl(
        \operatorname{Shv}_{\Sp}(X_{\mathrm{Zar}})
    \bigr)_{f^{-1}\mathcal O_Y/}.
\]
In particular, \(\mathcal O_X^\eta\) inherits a canonical structural
morphism
\[
    f^{-1}\mathcal O_Y
    \longrightarrow
    \mathcal O_X^\eta.
\]

By convention,
\[
    q_\eta^\sharp
    \colon
    \mathcal O_X^\eta
    \longrightarrow
    \mathcal O_X
\]
is the structure-sheaf morphism underlying the square-zero thickening.

There is a canonical fibre sequence of
\(\mathcal O_X^\eta\)-modules
\[
    \mathcal M
    \longrightarrow
    \mathcal O_X^\eta
    \xrightarrow{q_\eta^\sharp}
    \mathcal O_X,
\]
where \(\mathcal M\) and \(\mathcal O_X\) are regarded as
\(\mathcal O_X^\eta\)-modules by restriction of scalars along
\(q_\eta^\sharp\).
\end{construction}

\begin{theorem}[Square-zero thickenings of spectral schemes]
\label{thm:global-square-zero-spectral-scheme}
In \Ref{con:global-square-zero-extension}, suppose that
\[
    \mathcal M
    \in
    \QCoh(X)_{\geq0}.
\]
Then the spectrally ringed Zariski $\infty$-topos
\[
    X^\eta
    :=
    \bigl(
        X_{\mathrm{Zar}},
        \mathcal O_X^\eta
    \bigr)
\]
is a spectral scheme. The canonical morphism
\[
    i_\eta
    \colon
    X
    \hookrightarrow
    X^\eta
\]
is a closed immersion over $Y$, its underlying topological map is a
homeomorphism, and its affine restrictions are the structured square-zero
extensions of \Ref{def:structured-square-zero-extension}.
\end{theorem}
\begin{proof}
Because the pullback defining
\[
    \mathcal O_X^\eta
\]
is formed in commutative \(f^{-1}\mathcal O_Y\)-algebras,
\(\mathcal O_X^\eta\) carries a canonical structural morphism
\[
    f^{-1}\mathcal O_Y
    \longrightarrow
    \mathcal O_X^\eta.
\]
Thus, once the resulting ringed \(\infty\)-topos is shown to be a
spectral scheme, it is canonically a spectral scheme over \(Y\).

Let
\[
    U=\Spec(B)
    \longrightarrow
    X
\]
be an affine open. After shrinking an affine neighbourhood in \(Y\) if
necessary, suppose that the composite \(U\to Y\) factors through
\[
    V=\Spec(A)
    \longrightarrow
    Y.
\]
Write
\[
    M_U
    \in
    \Mod_B
\]
for the \(B\)-module corresponding to
\(\mathcal M|_U\).

By
\Ref{thm:global-cotangent-affine-compatibility}, restriction of
\(\eta\) to \(U\) is identified with a morphism
\[
    \eta_U
    \colon
    L_{B/A}
    \longrightarrow
    M_U[1].
\]
Restriction to \(U_{\mathrm{Zar}}\) preserves the specified pullback
defining \(\mathcal O_X^\eta\). Consequently,
\[
    \mathcal O_X^\eta|_U
\]
is canonically identified with the Zariski structure sheaf associated
to the commutative ring spectrum
\[
    B^{\eta_U}
    :=
    B
    \times_{B\oplus M_U[1]}
    B.
\]

Since \(B\) and \(M_U\) are connective,
\Ref{prop:fibre-structured-square-zero} shows that
\(B^{\eta_U}\) is connective. Hence the restriction of the ringed
\(\infty\)-topos
\[
    \bigl(
        X_{\mathrm{Zar}},
        \mathcal O_X^\eta
    \bigr)
\]
to the open subtopos \(U_{\mathrm{Zar}}\) is represented by the affine
spectral scheme
\[
    \Spec(B^{\eta_U}).
\]

These affine identifications are compatible on intersections. Indeed,
restriction to a smaller affine open is inverse image on sheaves, and
inverse image preserves the pullback defining the square-zero
extension. Equivalently, on coordinate algebras the compatibility is
the base-change compatibility of relative derivations and structured
square-zero extensions.

An affine open atlas of \(X\) therefore gives an affine open atlas of
\[
    X^\eta
    :=
    \bigl(
        X_{\mathrm{Zar}},
        \mathcal O_X^\eta
    \bigr).
\]
Thus \(X^\eta\) is a spectral scheme.

The morphism
\[
    q_\eta^\sharp
    \colon
    \mathcal O_X^\eta
    \longrightarrow
    \mathcal O_X
\]
defines a morphism of spectral schemes
\[
    i_\eta
    \colon
    X
    \longrightarrow
    X^\eta.
\]
Affine-locally it is represented by
\[
    B^{\eta_U}
    \longrightarrow
    B.
\]
By \Ref{prop:fibre-structured-square-zero}, this morphism is
surjective on \(\pi_0\) and its kernel is square-zero. The affine
closed-immersion criterion of
\Ref{def:affine-closed-immersion-spectral} therefore shows that
\(i_\eta\) is a closed immersion. Since closed immersions are
Zariski-local, the conclusion holds globally.

Moreover,
\[
    \Spec\pi_0(B)
    \longrightarrow
    \Spec\pi_0(B^{\eta_U})
\]
is defined by a square-zero ideal and is therefore a homeomorphism on
underlying topological spaces. These affine homeomorphisms glue, so
\[
    |X|
    \xrightarrow{\simeq}
    |X^\eta|
\]
is a homeomorphism.

Finally, the affine restriction of \(i_\eta\) over \(U=\Spec(B)\) is
exactly the structured square-zero extension classified by
\[
    \eta_U
    \colon
    L_{B/A}
    \longrightarrow
    M_U[1].
\]
Hence all assertions follow.
\end{proof}

\subsection{Intrinsic structured extensions}

\begin{definition}[Structured infinitesimal extension]
\label{def:structured-infinitesimal-extension}
Let
\[
    f
    \colon
    X
    \longrightarrow
    Y
\]
be a morphism of spectral schemes and let
\[
    \mathcal M
    \in
    \QCoh(X)_{\geq0}.
\]
A \textbf{structured infinitesimal extension of \(X\) over \(Y\) by
\(\mathcal M\)} is a morphism
\[
    \eta
    \colon
    L_{X/Y}
    \longrightarrow
    \mathcal M[1]
\]
together with the associated thickening
\[
    X
    \hookrightarrow
    X^\eta
\]
of \Ref{thm:global-square-zero-spectral-scheme}, including its specified
square-zero pullback identification. Its underlying infinitesimal
extension is the closed immersion obtained after forgetting the
classifying morphism and the specified pullback identification.

The resulting \(\infty\)-groupoid of structured infinitesimal extensions
of \(X\) over \(Y\) by \(\mathcal M\), with equivalences preserving the
classifying morphism and the specified pullback data, is denoted
\[
    \operatorname{SqZ}^{\mathrm{str}}_Y(X;\mathcal M).
\]
\end{definition}

\begin{theorem}[Global structured square-zero classification]
\label{thm:global-structured-square-zero-classification}
Let $f:X\to Y$ be a morphism of spectral schemes and let
$\mathcal M\in\QCoh(X)_{\geq0}$. Then there is a natural equivalence
\[
    \operatorname{SqZ}^{\mathrm{str}}_Y(X;\mathcal M)
    \simeq
    \Map_{\QCoh(X)}
    \bigl(
        L_{X/Y},
        \mathcal M[1]
    \bigr).
\]
\end{theorem}

\begin{proof}
The right side is the space of sheaf-level derivations. The construction
of \Ref{con:global-square-zero-extension} carries such a derivation to
its structured extension. For a fixed derivation, the space of pullback
models and identifications is contractible. The resulting sheaf-level
classification is developed in \cite[Section~17.1.3]{LurieSAG}.
\end{proof}

\section{Global deformation and obstruction theory}
\label{sec:global-spectral-deformation-theory}

\subsection{Deformations of spectral schemes}

\begin{definition}[Deformation over an infinitesimal extension]
\label{def:global-deformation-infinitesimal}
Let
\[
    i
    \colon
    S
    \hookrightarrow
    S'
\]
be a structured infinitesimal extension and let
\[
    f
    \colon
    X
    \longrightarrow
    S
\]
be a spectral scheme.

A \textbf{deformation of \(X\) over \(S'\)} is a pair
\[
    (X',\alpha)
\]
consisting of a spectral scheme
\[
    X'
    \longrightarrow
    S'
\]
and an equivalence
\[
    \alpha
    \colon
    X'\times_{S'}S
    \xrightarrow{\simeq}
    X
\]
over \(S\).

A morphism
\[
    (X'_1,\alpha_1)
    \longrightarrow
    (X'_2,\alpha_2)
\]
is an equivalence
\[
    u
    \colon
    X'_1
    \xrightarrow{\simeq}
    X'_2
\]
over \(S'\), together with the specified coherent compatibility
\[
    \alpha_2
    \circ
    \bigl(
        u\times_{S'}S
    \bigr)
    \simeq
    \alpha_1.
\]
The resulting \(\infty\)-groupoid is denoted
\[
    \operatorname{Def}_{S'/S}(X).
\]
\end{definition}

\begin{lemma}[Fixed-topos description of global deformations]
\label{lem:global-deformation-fixed-topos}
Let
\[
    i
    \colon
    S
    \hookrightarrow
    S'
\]
be a structured infinitesimal extension by
\[
    \mathcal I
    \in
    \QCoh(S)_{\geq0},
\]
classified by
\[
    \eta_S
    \colon
    L_S
    \longrightarrow
    \mathcal I[1].
\]
Let
\[
    f
    \colon
    X
    \longrightarrow
    S
\]
be a spectral scheme, put
\[
    \mathcal J
    :=
    f^*\mathcal I,
\]
and let
\[
    \eta_X
    :=
    f^*\eta_S
    \colon
    f^*L_S
    \longrightarrow
    \mathcal J[1].
\]
Then there is a natural equivalence
\[
    \operatorname{Def}_{S'/S}(X)
    \simeq
    \operatorname{fib}_{\eta_X}
    \left(
        \Map_{\QCoh(X)}
        \bigl(
            L_X,
            \mathcal J[1]
        \bigr)
        \longrightarrow
        \Map_{\QCoh(X)}
        \bigl(
            f^*L_S,
            \mathcal J[1]
        \bigr)
    \right),
\]
where the displayed morphism is induced by the canonical map
\[
    f^*L_S
    \longrightarrow
    L_X.
\]
Equivalently, a deformation is classified by a morphism
\[
    \widetilde\eta
    \colon
    L_X
    \longrightarrow
    \mathcal J[1]
\]
together with a specified homotopy
\[
    \widetilde\eta|_{f^*L_S}
    \simeq
    \eta_X.
\]
\end{lemma}
\begin{proof}
Let
\[
    (X',\alpha)
    \in
    \operatorname{Def}_{S'/S}(X).
\]
Write
\[
    f'
    \colon
    X'
    \longrightarrow
    S'
\]
for its structural morphism. The specified equivalence
\[
    \alpha
    \colon
    X'\times_{S'}S
    \xrightarrow{\simeq}
    X
\]
identifies the canonical morphism
\[
    j
    \colon
    X
    \longrightarrow
    X'
\]
with the base change of
\[
    i
    \colon
    S
    \longrightarrow
    S'.
\]

Since \(i\) is a square-zero closed immersion, so is \(j\). On
classical truncations, both are closed immersions defined by
square-zero ideals. Hence they induce homeomorphisms on underlying
topological spaces:
\[
    |S|
    \xrightarrow{\simeq}
    |S'|
    \qquad\text{and}\qquad
    |X|
    \xrightarrow{\simeq}
    |X'|.
\]
Spectral open subsets are detected on classical truncations. Therefore
pullback along \(i\) and \(j\) identifies the corresponding Zariski
open spectral subschemes and induces equivalences of Zariski
\(\infty\)-topoi
\[
    S_{\mathrm{Zar}}
    \xrightarrow{\simeq}
    S'_{\mathrm{Zar}}
\]
and
\[
    X_{\mathrm{Zar}}
    \xrightarrow{\simeq}
    X'_{\mathrm{Zar}}.
\]

Use these equivalences to transport the structure sheaves of \(S'\) and
\(X'\) to sheaves on \(S_{\mathrm{Zar}}\) and
\(X_{\mathrm{Zar}}\), respectively. The base extension becomes a
structured square-zero extension
\[
    \mathcal O_{S'}
    \longrightarrow
    \mathcal O_S
\]
on \(S_{\mathrm{Zar}}\), with coefficient module \(\mathcal I\) and
classifying morphism
\[
    \eta_S
    \colon
    L_S
    \longrightarrow
    \mathcal I[1].
\]

Similarly, the morphism \(j\) gives a structured square-zero extension
\[
    \mathcal O_{X'}
    \longrightarrow
    \mathcal O_X
\]
on \(X_{\mathrm{Zar}}\). Since \(j\) is obtained by base change from
\(i\), its coefficient module is canonically
\[
    \mathcal J
    =
    f^*\mathcal I.
\]

The sheaf-level tangent correspondence classifies the extension
\(\mathcal O_{X'}\to\mathcal O_X\) by a morphism
\[
    \widetilde\eta
    \colon
    L_X
    \longrightarrow
    \mathcal J[1].
\]
The morphism
\[
    f'
    \colon
    X'
    \longrightarrow
    S'
\]
and the specified identification of its reduction with
\(f:X\to S\) give a morphism from the source structured extension to
the base structured extension. Under functoriality of cotangent
complexes, this relative compatibility is precisely a specified
homotopy
\[
    \widetilde\eta|_{f^*L_S}
    \simeq
    f^*\eta_S
    =
    \eta_X.
\]
Thus a deformation determines a point of
\[
    \operatorname{fib}_{\eta_X}
    \left(
        \Map_{\QCoh(X)}
        \bigl(
            L_X,
            \mathcal J[1]
        \bigr)
        \longrightarrow
        \Map_{\QCoh(X)}
        \bigl(
            f^*L_S,
            \mathcal J[1]
        \bigr)
    \right).
\]

Conversely, suppose given a morphism
\[
    \widetilde\eta
    \colon
    L_X
    \longrightarrow
    \mathcal J[1]
\]
together with a specified homotopy
\[
    \widetilde\eta|_{f^*L_S}
    \simeq
    \eta_X.
\]

Use the equivalence
\[
    S_{\mathrm{Zar}}
    \simeq
    S'_{\mathrm{Zar}}
\]
to regard the geometric morphism
\[
    f
    \colon
    X_{\mathrm{Zar}}
    \longrightarrow
    S_{\mathrm{Zar}}
\]
as having target \(S'_{\mathrm{Zar}}\). Pulling the structured
extension
\[
    \mathcal O_{S'}
    \longrightarrow
    \mathcal O_S
\]
back to \(X_{\mathrm{Zar}}\) gives a structured square-zero extension
\[
    f^{-1}\mathcal O_{S'}
    \longrightarrow
    f^{-1}\mathcal O_S
\]
classified by
\[
    \eta_X
    \colon
    f^*L_S
    \longrightarrow
    \mathcal J[1].
\]

The specified homotopy
\[
    \widetilde\eta|_{f^*L_S}
    \simeq
    \eta_X
\]
is exactly the compatibility datum making the derivation classified by
\(\widetilde\eta\) a derivation over the structured base extension.
The relative sheaf-level tangent correspondence therefore constructs a
sheaf of connective commutative ring spectra
\[
    \mathcal O_{X'}
\]
on \(X_{\mathrm{Zar}}\), together with a commutative diagram
\[
\begin{tikzcd}[column sep=large, row sep=large]
    f^{-1}\mathcal O_{S'}
        \arrow[r]
        \arrow[d]
    &
    \mathcal O_{X'}
        \arrow[d]
    \\
    f^{-1}\mathcal O_S
        \arrow[r]
    &
    \mathcal O_X,
\end{tikzcd}
\]
in which the right vertical morphism is a structured square-zero
extension with coefficient module \(\mathcal J\), and the square is the
specified relative pullback square.

By
\Ref{thm:global-square-zero-spectral-scheme}, the spectrally ringed
Zariski \(\infty\)-topos
\[
    X'
    :=
    \bigl(
        X_{\mathrm{Zar}},
        \mathcal O_{X'}
    \bigr)
\]
is a spectral scheme. The upper horizontal morphism in the displayed
diagram supplies a morphism of spectrally ringed \(\infty\)-topoi
\[
    X'
    \longrightarrow
    S'.
\]
Base change along \(S\hookrightarrow S'\) recovers
\[
    \bigl(
        X_{\mathrm{Zar}},
        \mathcal O_X
    \bigr)
    =
    X,
\]
with its original structural morphism \(f:X\to S\). The specified
pullback identification gives an equivalence
\[
    X'\times_{S'}S
    \xrightarrow{\simeq}
    X.
\]
Thus the compatible derivation determines an object of
\[
    \operatorname{Def}_{S'/S}(X).
\]

The two constructions are inverse on objects. They are also inverse on
equivalences and all higher homotopies, because for a fixed derivation
the space of pullback models and specified pullback identifications is
contractible. Hence there is a natural equivalence
\[
    \operatorname{Def}_{S'/S}(X)
    \simeq
    \operatorname{fib}_{\eta_X}
    \left(
        \Map_{\QCoh(X)}
        \bigl(
            L_X,
            \mathcal J[1]
        \bigr)
        \longrightarrow
        \Map_{\QCoh(X)}
        \bigl(
            f^*L_S,
            \mathcal J[1]
        \bigr)
    \right).
\]
See
\cite[Sections~17.1.3 and~19.4.3]{LurieSAG}.
\end{proof}

\subsection{Split first-order deformations}

\begin{theorem}[Split first-order deformation space]
\label{thm:global-split-first-order-deformations}
Let
\[
    f
    \colon
    X
    \longrightarrow
    S
\]
be a morphism of spectral schemes and let
\[
    \mathcal I
    \in
    \QCoh(S)_{\geq0}.
\]
For the split square-zero thickening of \(S\) by \(\mathcal I\),
classified by the zero morphism
\[
    L_S
    \longrightarrow
    \mathcal I[1],
\]
the deformation \(\infty\)-groupoid of \(X\) is naturally equivalent to
\[
    \Map_{\QCoh(X)}
    \bigl(
        L_{X/S},
        f^*\mathcal I[1]
    \bigr).
\]
\end{theorem}

\begin{proof}
Put
\[
    \mathcal J
    :=
    f^*\mathcal I.
\]
Since the base extension is split, its classifying morphism is zero, and
hence the pulled-back classifying morphism
\[
    f^*L_S
    \longrightarrow
    \mathcal J[1]
\]
is also zero. By
\Ref{lem:global-deformation-fixed-topos},
the deformation \(\infty\)-groupoid is therefore the homotopy fibre over
zero of
\[
    \Map_{\QCoh(X)}
    \bigl(
        L_X,
        \mathcal J[1]
    \bigr)
    \longrightarrow
    \Map_{\QCoh(X)}
    \bigl(
        f^*L_S,
        \mathcal J[1]
    \bigr).
\]

The transitivity triangle
\[
    f^*L_S
    \longrightarrow
    L_X
    \longrightarrow
    L_{X/S}
\]
gives, after applying the exact contravariant mapping-spectrum functor
\[
    \map_{\QCoh(X)}(-,\mathcal J[1]),
\]
a fibre sequence of spectra
\[
    \map_{\QCoh(X)}
    \bigl(
        L_{X/S},
        \mathcal J[1]
    \bigr)
    \longrightarrow
    \map_{\QCoh(X)}
    \bigl(
        L_X,
        \mathcal J[1]
    \bigr)
    \longrightarrow
    \map_{\QCoh(X)}
    \bigl(
        f^*L_S,
        \mathcal J[1]
    \bigr).
\]
Passing to infinite loop spaces identifies the homotopy fibre over zero
with
\[
    \Map_{\QCoh(X)}
    \bigl(
        L_{X/S},
        \mathcal J[1]
    \bigr).
\]
Substituting
\[
    \mathcal J
    =
    f^*\mathcal I
\]
gives the asserted equivalence.
\end{proof}

\subsection{The Kodaira--Spencer obstruction}

\begin{theorem}[Global cotangent obstruction theory]
\label{thm:global-cotangent-obstruction}
Let
\[
    i
    \colon
    S
    \hookrightarrow
    S'
\]
be a structured infinitesimal extension by
\[
    \mathcal I
    \in
    \QCoh(S)_{\geq0},
\]
classified by
\[
    \eta_S
    \colon
    L_S
    \longrightarrow
    \mathcal I[1].
\]
Let
\[
    f
    \colon
    X
    \longrightarrow
    S
\]
be a spectral scheme and put
\[
    \mathcal J
    :=
    f^*\mathcal I.
\]
The transitivity triangle
\[
    f^*L_S
    \longrightarrow
    L_X
    \longrightarrow
    L_{X/S}
\]
determines a canonical Kodaira--Spencer obstruction
\[
    \operatorname{ob}(X/S')
    \in
    \pi_0
    \map_{\QCoh(X)}
    \bigl(
        L_{X/S},
        \mathcal J[2]
    \bigr).
\]
It satisfies:
\begin{enumerate}[label=(\roman*)]
    \item The obstruction vanishes if and only if \(X\) admits a
    deformation over \(S'\).

    \item When it vanishes, \(\operatorname{Def}_{S'/S}(X)\) is a torsor
    under
    \[
        \Omega^\infty
        \map_{\QCoh(X)}
        \bigl(
            L_{X/S},
            \mathcal J[1]
        \bigr).
    \]

    \item Let
    \[
        \xi
        \in
        \operatorname{Def}_{S'/S}(X)
    \]
    be a fixed deformation. Write
    \[
        \operatorname{Aut}(\xi)
        :=
        \Omega_\xi
        \operatorname{Def}_{S'/S}(X)
    \]
    for its automorphism space in the deformation $\infty$-groupoid.
    Thus its elements are automorphisms over \(S'\) which preserve the
    specified identification of the reduction with \(X\), together with
    the required coherent homotopy. There is a natural equivalence
    \[
        \operatorname{Aut}(\xi)
        \simeq
        \Omega^\infty
        \map_{\QCoh(X)}
        \bigl(
            L_{X/S},
            \mathcal J
        \bigr).
    \]
\end{enumerate}
\end{theorem}

\begin{proof}
Pulling back the classifying morphism of the base extension gives
\[
    \eta_X
    :=
    f^*\eta_S
    \colon
    f^*L_S
    \longrightarrow
    \mathcal J[1].
\]
By
\Ref{lem:global-deformation-fixed-topos},
there is a natural equivalence
\[
    \operatorname{Def}_{S'/S}(X)
    \simeq
    \operatorname{fib}_{\eta_X}
    \left(
        \Map_{\QCoh(X)}
        \bigl(
            L_X,
            \mathcal J[1]
        \bigr)
        \longrightarrow
        \Map_{\QCoh(X)}
        \bigl(
            f^*L_S,
            \mathcal J[1]
        \bigr)
    \right).
\]

Applying the exact contravariant functor
\[
    \map_{\QCoh(X)}(-,\mathcal J[1])
\]
to the transitivity triangle
\[
    f^*L_S
    \longrightarrow
    L_X
    \longrightarrow
    L_{X/S}
\]
gives a fibre sequence of spectra
\[
\begin{aligned}
    \map_{\QCoh(X)}
    \bigl(
        L_{X/S},
        \mathcal J[1]
    \bigr)
    &\longrightarrow
    \map_{\QCoh(X)}
    \bigl(
        L_X,
        \mathcal J[1]
    \bigr)
    \\
    &\longrightarrow
    \map_{\QCoh(X)}
    \bigl(
        f^*L_S,
        \mathcal J[1]
    \bigr)
    \\
    &\xrightarrow{\partial}
    \map_{\QCoh(X)}
    \bigl(
        L_{X/S},
        \mathcal J[2]
    \bigr).
\end{aligned}
\]
Define
\[
    \operatorname{ob}(X/S')
    :=
    \partial(\eta_X).
\]

The point \(\eta_X\) lifts through
\[
    \Map_{\QCoh(X)}
    \bigl(
        L_X,
        \mathcal J[1]
    \bigr)
    \longrightarrow
    \Map_{\QCoh(X)}
    \bigl(
        f^*L_S,
        \mathcal J[1]
    \bigr)
\]
if and only if its boundary is null. By the homotopy-fibre description,
such a lift is exactly a deformation of \(X\) over \(S'\). This proves
\textup{(i)}.

When the homotopy fibre is nonempty, translation by a chosen point
identifies it as a torsor under the homotopy fibre over zero. By the
displayed fibre sequence, that homotopy fibre is
\[
    \Omega^\infty
    \map_{\QCoh(X)}
    \bigl(
        L_{X/S},
        \mathcal J[1]
    \bigr).
\]
This proves \textup{(ii)}.

For a fixed deformation \(\xi\), its automorphism space in the
deformation $\infty$-groupoid is the based loop space
\[
    \Omega_\xi
    \operatorname{Def}_{S'/S}(X).
\]
Using the torsor description gives
\[
\begin{aligned}
    \operatorname{Aut}(\xi)
    &\simeq
    \Omega
    \Omega^\infty
    \map_{\QCoh(X)}
    \bigl(
        L_{X/S},
        \mathcal J[1]
    \bigr)
    \\
    &\simeq
    \Omega^\infty
    \map_{\QCoh(X)}
    \bigl(
        L_{X/S},
        \mathcal J
    \bigr).
\end{aligned}
\]
This proves \textup{(iii)}. See
\cite[Section~19.4.3]{LurieSAG}.
\end{proof}

\begin{corollary}[Affine compatibility of global obstruction theory]
\label{cor:global-obstruction-affine-compatibility}
Suppose
\[
    S=\Spec(A),
    \qquad
    S'=\Spec(A'),
    \qquad
    X=\Spec(B).
\]
Then every deformation of \(X\) over \(S'\) is affine, and under the
resulting identification of deformation \(\infty\)-groupoids,
\Ref{thm:global-cotangent-obstruction}
identifies canonically with
\Ref{thm:affine-cotangent-obstruction}.
\end{corollary}

\begin{proof}
Let
\[
    (X',\alpha)
    \in
    \operatorname{Def}_{S'/S}(X).
\]
The morphism
\[
    X
    \simeq
    X'\times_{S'}S
    \longrightarrow
    X'
\]
is the base change of the square-zero closed immersion
\[
    S
    \longrightarrow
    S'.
\]
It is therefore a square-zero closed immersion and induces a
homeomorphism on underlying topological spaces. On classical
truncations,
\[
    X_{\mathrm{cl}}
    \longrightarrow
    X'_{\mathrm{cl}}
\]
is a nilpotent closed immersion.

Since \(X_{\mathrm{cl}}\) is affine, the nilpotent thickening
\(X'_{\mathrm{cl}}\) is affine. Affineness of spectral schemes is
detected on classical truncations, so \(X'\) is affine. Write
\[
    X'
    \simeq
    \Spec(B')
\]
for a connective commutative \(A'\)-algebra \(B'\). The specified
equivalence
\[
    X'\times_{S'}S
    \simeq
    X
\]
then identifies contravariantly with an equivalence
\[
    B'\otimes_{A'}A
    \xrightarrow{\simeq}
    B.
\]
Thus the global and affine definitions of deformation determine the same
\(\infty\)-groupoid.

Under the affine comparison theorem,
\[
    L_X
    \simeq
    L_B,
    \qquad
    f^*L_S
    \simeq
    L_A\otimes_AB,
    \qquad
    L_{X/S}
    \simeq
    L_{B/A},
\]
and
\[
    f^*\mathcal I
    \simeq
    B\otimes_AI.
\]
Moreover, quasi-coherent mapping spectra on \(\Spec(B)\) identify with
mapping spectra in \(\Mod_B\). Consequently, the global transitivity
triangle, its boundary morphism, the obstruction class, the deformation
torsor, and the automorphism space identify respectively with those of
\Ref{thm:affine-cotangent-obstruction}.
\end{proof}

\section{Quasi-smooth closed immersions and derived zero loci}
\label{sec:quasi-smooth-closed-immersions}

\subsection{Conormal modules and quasi-smoothness}

\begin{definition}[Conormal module]
\label{def:conormal-module}
Let
\[
    i
    \colon
    Z
    \hookrightarrow
    X
\]
be a closed immersion of spectral schemes. Its \textbf{conormal module}
is
\[
    N^*_{Z/X}
    :=
    L_{Z/X}[-1]
    \in
    \QCoh(Z).
\]
\end{definition}

\begin{definition}[Quasi-smooth or regular closed immersion]
\label{def:quasi-smooth-closed-immersion}
A closed immersion
\[
    i
    \colon
    Z
    \hookrightarrow
    X
\]
is \textbf{quasi-smooth} if:
\begin{enumerate}[label=(\roman*)]
    \item it is locally of finite presentation;

    \item the conormal module
    \[
        N^*_{Z/X}
        =
        L_{Z/X}[-1]
    \]
    is locally free of finite rank.
\end{enumerate}
A quasi-smooth closed immersion will also be called a
\textbf{regular closed immersion} in this work.
\end{definition}

\begin{definition}[Normal module and normal bundle]
\label{def:normal-bundle}
For a quasi-smooth closed immersion
\[
    i
    \colon
    Z
    \hookrightarrow
    X,
\]
define the \textbf{normal module}
\[
    N_{Z/X}
    :=
    (N^*_{Z/X})^\vee.
\]
Define the \textbf{normal bundle} by
\[
    \mathbf N_{Z/X}
    :=
    \mathbb V_Z(N^*_{Z/X}).
\]
Since \(N^*_{Z/X}\) is finite locally free, finite-projective biduality
gives a canonical equivalence
\[
    (N_{Z/X})^\vee
    =
    \bigl((N^*_{Z/X})^\vee\bigr)^\vee
    \simeq
    N^*_{Z/X}.
\]
Consequently, there is a canonical equivalence
\[
    \mathbf N_{Z/X}
    \simeq
    \mathbf V_Z(N_{Z/X}).
\]
\end{definition}

\begin{proposition}[Base change of quasi-smooth closed immersions]
\label{prop:quasi-smooth-closed-base-change}
Quasi-smooth closed immersions are stable under arbitrary base change.
For every morphism $Y\to X$, there is a natural equivalence
\[
    N^*_{Z\times_XY/Y}
    \simeq
    N^*_{Z/X}|_{Z\times_XY}.
\]
\end{proposition}

\begin{proof}
Closed immersions and morphisms locally of finite presentation are stable
under base change by the results of
\Ref{chap:spectral-algebraic-geometry}. Global cotangent base change gives
\[
    L_{Z\times_XY/Y}
    \simeq
    L_{Z/X}|_{Z\times_XY}.
\]
Shift by $[-1]$ and apply
\Ref{prop:base-change-locally-free}.
\end{proof}

\subsection{Derived zero loci}

\begin{construction}[Zero locus of a section]
\label{con:derived-zero-locus}
Let $\mathcal E$ be a finite locally free quasi-coherent module on a
spectral scheme $X$, and let
\[
    \sigma
    \colon
    X
    \longrightarrow
    \mathbb V_X(\mathcal E)
\]
be a section. Under
\Ref{prop:linear-space-functor-points}, the section $\sigma$ corresponds
to a morphism of quasi-coherent $\mathcal O_X$-modules
\[
    s
    \colon
    \mathcal E
    \longrightarrow
    \mathcal O_X.
\]
Equivalently, finite-projective duality identifies $s$ with a section
\[
    s^\vee
    \colon
    \mathcal O_X
    \longrightarrow
    \mathcal E^\vee.
\]
Thus this construction is the usual derived zero locus of a section of
$\mathcal E^\vee$; the use of $\mathcal E\to\mathcal O_X$ reflects
the convention
\[
    \mathbb V_X(\mathcal E)
    =
    \Spec_X\operatorname{Sym}_{\mathcal O_X}(\mathcal E).
\]

The section $\sigma$ and the zero section
\[
    0
    \colon
    X
    \longrightarrow
    \mathbb V_X(\mathcal E)
\]
determine a derived fibre product
\[
    Z(s)
    :=
    X
    \times_{\mathbb V_X(\mathcal E)}
    X,
\]
where the two morphisms to $\mathbb V_X(\mathcal E)$ are respectively
$\sigma$ and the zero section. It is called the
\textbf{derived zero locus of $s$}.
\end{construction}

\begin{proposition}[Derived zero loci are quasi-smooth]
\label{prop:derived-zero-locus-quasi-smooth}
In \Ref{con:derived-zero-locus}, the projection
\[
    Z(s)
    \longrightarrow
    X
\]
is a quasi-smooth closed immersion. Its conormal module is canonically
\[
    N^*_{Z(s)/X}
    \simeq
    \mathcal E|_{Z(s)}.
\]
\end{proposition}

\begin{proof}
The assertion is local on \(X\). After a Zariski localization, write
\[
    X
    =
    \Spec(A)
\]
and choose an equivalence
\[
    \mathcal E
    \simeq
    A^{\oplus n}.
\]
Put
\[
    R
    :=
    \operatorname{Sym}_A(A^{\oplus n}).
\]

The zero section and the section \(\sigma\) of
\Ref{con:derived-zero-locus} correspond to two morphisms of commutative
\(A\)-algebras
\[
    \epsilon_0,\epsilon_s
    \colon
    R
    \rightrightarrows
    A,
\]
where \(\epsilon_0\) is the zero augmentation.

To distinguish the two \(R\)-algebra structures, write \(A_s\) for the
copy of \(A\) equipped with the structural morphism \(\epsilon_s\), and
write \(A_0\) for the copy of \(A\) equipped with the structural
morphism \(\epsilon_0\). The coordinate algebra of the derived zero
locus is the pushout
\[
    C
    :=
    A_s
    \otimes_R
    A_0.
\]
We regard \(C\) as an \(A=A_s\)-algebra through the canonical morphism
\[
    A_s
    \longrightarrow
    C.
\]

We first analyze the zero augmentation
\[
    \epsilon_0
    \colon
    R
    \longrightarrow
    A_0.
\]
Consider the composable morphisms
\[
    A
    \longrightarrow
    R
    \xrightarrow{\epsilon_0}
    A_0.
\]
Under the canonical identification \(A_0\simeq A\) of underlying
commutative \(A\)-algebras, their composite is the identity of \(A\).
The transitivity triangle therefore gives a cofibre sequence
\[
    L_{R/A}\otimes_RA_0
    \longrightarrow
    L_{A_0/A}
    \longrightarrow
    L_{A_0/R}.
\]
Since
\[
    L_{A_0/A}
    \simeq
    0
\]
and, by the free-algebra calculation,
\[
    L_{R/A}
    \simeq
    R\otimes_AA^{\oplus n},
\]
there is a natural equivalence
\[
    L_{A_0/R}
    \simeq
    A_0^{\oplus n}[1].
\]
In particular, \(L_{A_0/R}\) is a perfect \(A_0\)-module.

On degree-zero homotopy, the zero augmentation identifies
\[
    \pi_0(A_0)
    \simeq
    \pi_0(R)/(t_1,\ldots,t_n),
\]
so \(\pi_0(A_0)\) is finitely presented as a
\(\pi_0(R)\)-algebra. It follows from
\Ref{thm:cotangent-finite-presentation} that
\[
    R
    \longrightarrow
    A_0
\]
is of finite presentation. Since it is surjective on \(\pi_0\), it
represents a closed immersion
\[
    \Spec(A_0)
    \longrightarrow
    \Spec(R).
\]

The projection
\[
    \Spec(C)
    \longrightarrow
    \Spec(A_s)
\]
is the base change of this zero section along the morphism
\[
    \Spec(A_s)
    \longrightarrow
    \Spec(R)
\]
represented by \(\sigma\). Hence it is a closed immersion of finite
presentation.

Applying cotangent base change to the pushout square
\[
\begin{tikzcd}[column sep=large, row sep=large]
    R
        \arrow[r, "\epsilon_0"]
        \arrow[d, "\epsilon_s"']
    &
    A_0
        \arrow[d]
    \\
    A_s
        \arrow[r]
    &
    C
\end{tikzcd}
\]
gives
\[
    L_{C/A_s}
    \simeq
    L_{A_0/R}\otimes_{A_0}C
    \simeq
    C^{\oplus n}[1].
\]
Consequently,
\[
    L_{C/A_s}[-1]
    \simeq
    C^{\oplus n}.
\]
Thus
\[
    \Spec(C)
    \longrightarrow
    \Spec(A_s)
\]
is quasi-smooth, and its conormal module is the restriction of
\(A^{\oplus n}\).

These affine identifications are natural under further Zariski
localization and therefore glue to a canonical equivalence
\[
    N^*_{Z(s)/X}
    \simeq
    \mathcal E|_{Z(s)}.
\]
\end{proof}

\begin{lemma}[Universal equations and the conormal basis]
\label{lem:universal-equations-conormal-basis}
Let \(A\) be a connective commutative ring spectrum and put
\[
    R
    :=
    \operatorname{Sym}_A(A^{\oplus n}).
\]
Let
\[
    \epsilon_0,
    \epsilon_f
    \colon
    R
    \rightrightarrows
    A
\]
be respectively the zero augmentation and the morphism classified by
points
\[
    f_1,\ldots,f_n
    \in
    \Omega^\infty A.
\]
To distinguish the two factors, write \(A_f\) for the target of
\(\epsilon_f\) and \(A_0\) for the target of \(\epsilon_0\), and put
\[
    C
    :=
    A_f\otimes_RA_0.
\]
Regard \(C\) as an \(A\)-algebra through
\[
    A=A_f
    \longrightarrow
    C
\]
and put
\[
    K_C
    :=
    \operatorname{fib}(A\to C).
\]

For every \(1\leq i\leq n\), the \(i\)-th universal equation, together
with its canonical nullhomotopy in \(C\), determines a class
\[
    \lambda_i
    \in
    \pi_0(K_C).
\]
Under the canonical equivalence
\[
    L_{C/A}[-1]
    \simeq
    C^{\oplus n},
\]
the derived-conormal comparison
\[
    C\otimes_AK_C
    \longrightarrow
    L_{C/A}[-1]
\]
carries the image of \(\lambda_i\) to the \(i\)-th standard basis vector
of
\[
    \pi_0(C)^{\oplus n}.
\]
\end{lemma}

\begin{proof}
Let
\[
    x_1,\ldots,x_n
\]
denote the universal generators of
\[
    R
    =
    \operatorname{Sym}_A(A^{\oplus n}).
\]
The universal derivation gives the canonical equivalence
\[
    L_{R/A}
    \simeq
    R\otimes_AA^{\oplus n},
\]
under which the differential of \(x_i\) is the \(i\)-th standard basis
vector.

Apply transitivity to
\[
    A
    \longrightarrow
    R
    \xrightarrow{\epsilon_0}
    A_0.
\]
Since the composite is the identity and
\[
    L_{A_0/A}
    \simeq
    0,
\]
the transitivity triangle gives a canonical equivalence
\[
    L_{A_0/R}[-1]
    \simeq
    L_{R/A}\otimes_RA_0
    \simeq
    A_0^{\oplus n}.
\]
Under this equivalence, the class determined by \(x_i\) and its
nullhomotopy under the zero augmentation maps to the \(i\)-th standard
basis vector. This is the boundary description of the transitivity
morphism applied to the universal derivation.

The pushout square
\[
\begin{tikzcd}[column sep=large, row sep=large]
    R
        \arrow[r, "\epsilon_0"]
        \arrow[d, "\epsilon_f"']
    &
    A_0
        \arrow[d]
    \\
    A_f
        \arrow[r]
    &
    C
\end{tikzcd}
\]
and cotangent base change give
\[
    L_{C/A_f}
    \simeq
    L_{A_0/R}\otimes_{A_0}C.
\]
Consequently,
\[
    L_{C/A}[-1]
    \simeq
    C^{\oplus n}.
\]

The canonical homotopy between the images of \(x_i\) through the two
legs of the pushout determines the class
\[
    \lambda_i
    \in
    \pi_0(K_C).
\]
Naturality of the relative-cofibre-to-cotangent morphism under the
pushout square identifies the image of \(\lambda_i\) with the base change
of the corresponding universal class for
\(\epsilon_0:R\to A_0\). That universal class maps to the \(i\)-th
standard basis vector by the preceding calculation. Its base change
therefore maps to the \(i\)-th standard basis vector of
\[
    \pi_0(C)^{\oplus n}.
\]
\end{proof}

\subsection{Local zero-locus structure}

\begin{theorem}[Local zero-locus characterization]
\label{thm:quasi-smooth-local-zero-locus}
Let
\[
    i
    \colon
    Z
    \hookrightarrow
    X
\]
be a closed immersion of spectral schemes. The following conditions are
equivalent.
\begin{enumerate}[label=(\roman*)]
    \item The immersion \(i\) is quasi-smooth.

    \item Zariski-locally on \(X\), the immersion is equivalent to the
    derived zero locus of a section of a finite locally free
    quasi-coherent module, or equivalently, in the
    \(\mathbb V\)-convention, of a morphism
    \[
        \mathcal E
        \longrightarrow
        \mathcal O_X
    \]
    with \(\mathcal E\) finite locally free.

    \item Zariski-locally on \(X\), there is an open spectral subscheme
    \[
        U
        \subseteq
        X,
    \]
    an integer \(n\geq0\), and a section
    \[
        s
        \colon
        U
        \longrightarrow
        \mathbb A_U^n
    \]
    such that
    \[
        Z\times_XU
        \longrightarrow
        U
    \]
    is equivalent to the base change of the zero section
    \[
        0
        \colon
        U
        \longrightarrow
        \mathbb A_U^n
    \]
    along \(s\).
\end{enumerate}
\end{theorem}

\begin{proof}
The equivalence of \textup{(ii)} and \textup{(iii)} follows by
Zariski-locally trivializing the finite locally free module. The
implication
\[
    \textup{(ii)}
    \Longrightarrow
    \textup{(i)}
\]
is \Ref{prop:derived-zero-locus-quasi-smooth}. It remains to prove
\[
    \textup{(i)}
    \Longrightarrow
    \textup{(ii)}.
\]

The assertion is local on \(X\). On the open complement of \(Z\), the
immersion is empty. The empty immersion is the derived zero locus of the
unit section
\[
    \mathcal O_X
    \longrightarrow
    \mathcal O_X.
\]
Indeed, affine-locally this zero locus is represented by
\[
    A
    \otimes_{\operatorname{Sym}_A(A)}
    A,
\]
where the universal generator is sent to \(1\) on one factor and to
\(0\) on the other. Its degree-zero homotopy ring is
\[
    \pi_0(A)/(1)
    =
    0.
\]
Since the resulting commutative ring spectrum is connective, it is the
zero ring spectrum, and its spectrum is empty.
It is therefore enough to work in a neighbourhood of a point of \(Z\).

Choose an affine neighbourhood of the selected point and write the
immersion as
\[
    \Spec(B)
    \longrightarrow
    \Spec(A),
\]
where
\[
    A
    \longrightarrow
    B
\]
is a morphism of connective commutative ring spectra which is surjective
on \(\pi_0\) and of finite presentation.

The conormal module
\[
    L_{B/A}[-1]
\]
is locally free of finite rank on \(\Spec(B)\). Choose an open
neighbourhood of the selected point in \(\Spec(B)\) on which it is free
of rank \(n\). Since
\[
    \Spec(B)
    \longrightarrow
    \Spec(A)
\]
is a closed immersion, this open neighbourhood is the intersection of
\(\Spec(B)\) with an open neighbourhood of the selected point in
\(\Spec(A)\). Spectral open subsets are detected on classical
truncations. After shrinking the ambient affine \(\Spec(A)\), we may
therefore choose an equivalence
\[
    L_{B/A}[-1]
    \simeq
    B^{\oplus n}.
\]

Put
\[
    K
    :=
    \operatorname{fib}(A\to B),
\]
and write
\[
    A_0
    :=
    \pi_0(A),
    \qquad
    B_0
    :=
    \pi_0(B),
    \qquad
    I
    :=
    \ker(A_0\to B_0).
\]
Since \(A_0\to B_0\) is surjective, there is an identification
\[
    B_0
    \simeq
    A_0/I.
\]

By \Ref{thm:derived-conormal-comparison}, the canonical morphism
\[
    B\otimes_AK
    \longrightarrow
    L_{B/A}[-1]
\]
induces an isomorphism
\[
    \pi_0(B\otimes_AK)
    \xrightarrow{\simeq}
    \pi_0
    \bigl(
        L_{B/A}[-1]
    \bigr).
\]
Because \(B\) and \(K\) are connective, the connective tensor-product
spectral sequence gives
\[
    \pi_0(B\otimes_AK)
    \simeq
    B_0\otimes_{A_0}\pi_0(K).
\]
Using
\[
    B_0
    \simeq
    A_0/I,
\]
this becomes
\[
    \pi_0(B\otimes_AK)
    \simeq
    \pi_0(K)/I\pi_0(K).
\]
Consequently, the natural morphism
\[
    \pi_0(K)
    \longrightarrow
    \pi_0(B\otimes_AK)
\]
is surjective.

Let
\[
    e_1,\ldots,e_n
\]
be the basis of
\[
    \pi_0
    \bigl(
        L_{B/A}[-1]
    \bigr)
    \simeq
    B_0^{\oplus n}
\]
determined by the chosen trivialization. Choose points
\[
    \widetilde\kappa_1,\ldots,\widetilde\kappa_n
    \in
    \Omega^\infty K
\]
whose connected components map under
\[
    \pi_0(K)
    \longrightarrow
    \pi_0(B\otimes_AK)
    \xrightarrow{\simeq}
    \pi_0
    \bigl(
        L_{B/A}[-1]
    \bigr)
\]
to
\[
    e_1,\ldots,e_n.
\]
Write
\[
    \kappa_i
    :=
    [\widetilde\kappa_i]
    \in
    \pi_0(K).
\]

Let
\[
    \widetilde f_i
    \in
    \Omega^\infty A
\]
be the image of \(\widetilde\kappa_i\), and put
\[
    f_i
    :=
    [\widetilde f_i]
    \in
    A_0.
\]
Since \(\widetilde\kappa_i\) is a point of the homotopy fibre of
\[
    \Omega^\infty A
    \longrightarrow
    \Omega^\infty B
\]
over zero, it includes a specified path from the image of
\(\widetilde f_i\) in \(\Omega^\infty B\) to zero. In particular,
\[
    f_i
    \in
    I.
\]

The surjection of
\Ref{lem:derived-conormal-classical-quotient}
\[
    \pi_0(B\otimes_AK)
    \longrightarrow
    I/I^2
\]
carries the class represented by \(\kappa_i\) to the class of \(f_i\).
Since the classes represented by the \(\kappa_i\) form a basis of
\(\pi_0(B\otimes_AK)\), the classes of
\[
    f_1,\ldots,f_n
\]
generate \(I/I^2\).

We next verify that \(I\) is finitely generated. Since \(A\to B\) is of
finite presentation, \(B\) is compact in
\[
    \CAlg(\Sp)_{A/}.
\]
We claim that \(B_0\) is compact in
\[
    \operatorname{CRing}_{A_0/}.
\]
Let
\[
    \{T_\alpha\}_\alpha
\]
be a filtered diagram of ordinary commutative \(A_0\)-algebras. The
Eilenberg--Mac Lane construction carries it to a filtered diagram of
connective commutative \(A\)-algebras, where the \(A\)-algebra structure
is induced by
\[
    A
    \longrightarrow
    HA_0
    \longrightarrow
    HT_\alpha.
\]
Filtered colimits of discrete connective commutative ring spectra remain
discrete, so there is a natural equivalence
\[
    H
    \left(
        \varinjlim_\alpha T_\alpha
    \right)
    \simeq
    \varinjlim_\alpha HT_\alpha.
\]
Using the adjunction
\[
    \pi_0
    \dashv
    H
\]
and compactness of \(B\), we obtain
\[
\begin{aligned}
    \Map_{\operatorname{CRing}_{A_0/}}
    \left(
        B_0,
        \varinjlim_\alpha T_\alpha
    \right)
    &\simeq
    \Map_{\CAlg(\Sp)_{A/}}
    \left(
        B,
        H
        \left(
            \varinjlim_\alpha T_\alpha
        \right)
    \right)
    \\
    &\simeq
    \Map_{\CAlg(\Sp)_{A/}}
    \left(
        B,
        \varinjlim_\alpha HT_\alpha
    \right)
    \\
    &\simeq
    \varinjlim_\alpha
    \Map_{\CAlg(\Sp)_{A/}}
    \left(
        B,
        HT_\alpha
    \right)
    \\
    &\simeq
    \varinjlim_\alpha
    \Map_{\operatorname{CRing}_{A_0/}}
    \left(
        B_0,
        T_\alpha
    \right).
\end{aligned}
\]
Thus \(B_0\) is finitely presented as an \(A_0\)-algebra.

Since
\[
    A_0
    \longrightarrow
    B_0
\]
is surjective and finitely presented, its kernel \(I\) is finitely
generated. Put
\[
    Q
    :=
    I/(f_1,\ldots,f_n).
\]
The module \(Q\) is finitely generated. Since the classes of the \(f_i\)
generate \(I/I^2\), one has
\[
    I
    =
    (f_1,\ldots,f_n)+I^2,
\]
and hence
\[
    Q
    =
    IQ.
\]

Let the selected point of \(Z\) correspond to a prime ideal
\[
    \mathfrak p
    \subseteq
    A_0
\]
containing \(I\). Localizing at \(\mathfrak p\), the ideal
\(I_{\mathfrak p}\) lies in the maximal ideal of
\((A_0)_{\mathfrak p}\), and
\[
    Q_{\mathfrak p}
    =
    I_{\mathfrak p}Q_{\mathfrak p}.
\]
Nakayama's lemma gives
\[
    Q_{\mathfrak p}
    =
    0.
\]
Since \(Q\) is finitely generated, after shrinking the ambient affine
open around the selected point, one has
\[
    Q
    =
    0.
\]
Thus, after this localization,
\[
    I
    =
    (f_1,\ldots,f_n).
\]

Put
\[
    R
    :=
    \operatorname{Sym}_A(A^{\oplus n}).
\]
Let
\[
    \epsilon_0
    \colon
    R
    \longrightarrow
    A
\]
be the zero augmentation. Let
\[
    u_f
    \colon
    A^{\oplus n}
    \longrightarrow
    A
\]
be the \(A\)-module morphism whose components correspond to the points
\[
    \widetilde f_1,\ldots,\widetilde f_n
    \in
    \Omega^\infty A.
\]
By the free-algebra adjunction, \(u_f\) determines a morphism of
commutative \(A\)-algebras
\[
    \epsilon_f
    \colon
    R
    \longrightarrow
    A.
\]
Define
\[
    C
    :=
    A\otimes_RA,
\]
where the two morphisms from \(R\) are \(\epsilon_f\) and
\(\epsilon_0\), and regard \(C\) as an \(A\)-algebra through the factor
corresponding to \(\epsilon_f\).

After composing with
\[
    A
    \longrightarrow
    B,
\]
the maps \(\epsilon_0\) and \(\epsilon_f\) correspond under the
free-algebra adjunction to the zero map
\[
    A^{\oplus n}
    \longrightarrow
    B
\]
and the map whose components are the images of
\(\widetilde f_1,\ldots,\widetilde f_n\). The specified paths carried by
\(\widetilde\kappa_1,\ldots,\widetilde\kappa_n\) give a homotopy between
these two module maps and hence a homotopy between the corresponding
commutative \(A\)-algebra maps
\[
    R
    \rightrightarrows
    B.
\]
By the pushout universal property, this determines an \(A\)-algebra
morphism
\[
    C
    \longrightarrow
    B.
\]

Classical truncation preserves this pushout. Therefore
\[
    \pi_0(C)
    \simeq
    A_0
    \otimes_{\pi_0(R)}
    A_0.
\]
Under the two maps from \(\pi_0(R)\), the universal variables are sent
respectively to
\[
    f_1,\ldots,f_n
    \qquad\text{and}\qquad
    0.
\]
Consequently,
\[
    \pi_0(C)
    \simeq
    A_0/(f_1,\ldots,f_n)
    \simeq
    A_0/I
    \simeq
    B_0.
\]
Thus
\[
    \pi_0(C)
    \xrightarrow{\simeq}
    \pi_0(B).
\]

By \Ref{prop:derived-zero-locus-quasi-smooth}, there is a natural
equivalence
\[
    L_{C/A}[-1]
    \simeq
    C^{\oplus n}.
\]
The transitivity triangle for
\[
    A
    \longrightarrow
    C
    \longrightarrow
    B
\]
contains the morphism
\[
    L_{C/A}\otimes_CB
    \longrightarrow
    L_{B/A}.
\]
After shifting by \([-1]\), this gives
\[
    \bigl(
        L_{C/A}\otimes_CB
    \bigr)[-1]
    \longrightarrow
    L_{B/A}[-1].
\]

Put
\[
    K_C
    :=
    \operatorname{fib}(A\to C),
    \qquad
    K_B
    :=
    \operatorname{fib}(A\to B)=K.
\]
The morphism \(C\to B\) induces a morphism
\[
    K_C
    \longrightarrow
    K_B.
\]
Since both \(A\to C\) and \(A\to B\) are surjective on \(\pi_0\),
naturality in
\Ref{thm:derived-conormal-comparison}
gives a commutative square of \(B\)-modules
\[
\begin{tikzcd}[column sep=large, row sep=large]
    B\otimes_AK_C
        \arrow[r]
        \arrow[d]
    &
    \bigl(
        L_{C/A}\otimes_CB
    \bigr)[-1]
        \arrow[d]
    \\
    B\otimes_AK_B
        \arrow[r]
    &
    L_{B/A}[-1].
\end{tikzcd}
\]

For each \(i\), the \(i\)-th universal equation of \(C\), together with
its canonical nullhomotopy on \(\Spec(C)\), determines a class
\[
    \lambda_i
    \in
    \pi_0(K_C).
\]
The morphism
\[
    C
    \longrightarrow
    B
\]
was constructed using the specified nullhomotopy carried by
\(\widetilde\kappa_i\). Consequently, the image of \(\lambda_i\) under
\[
    \pi_0(K_C)
    \longrightarrow
    \pi_0(K_B)
\]
is \(\kappa_i\).

By
\Ref{lem:universal-equations-conormal-basis},
the upper horizontal derived-conormal comparison, followed by the
equivalence
\[
    \bigl(
        L_{C/A}\otimes_CB
    \bigr)[-1]
    \simeq
    B^{\oplus n},
\]
carries the image of \(\lambda_i\) to the \(i\)-th standard basis
vector. By the choice of \(\kappa_i\), the lower horizontal comparison
carries its image to
\[
    e_i
    \in
    \pi_0
    \bigl(
        L_{B/A}[-1]
    \bigr).
\]
Commutativity of the square therefore shows that the morphism
\[
    \bigl(
        L_{C/A}\otimes_CB
    \bigr)[-1]
    \longrightarrow
    L_{B/A}[-1]
\]
carries the \(i\)-th standard basis vector to \(e_i\). Hence it induces
an isomorphism on \(\pi_0\).

Both its source and target are finite projective \(B\)-modules. For a
finite projective \(B\)-module \(P\), the flat homotopy-group formula
gives
\[
    \pi_k(P)
    \simeq
    \pi_k(B)
    \otimes_{B_0}
    \pi_0(P)
\]
for every integer \(k\). It follows that a morphism between finite
projective \(B\)-modules which is an isomorphism on \(\pi_0\) is an
equivalence. Therefore
\[
    L_{C/A}\otimes_CB
    \xrightarrow{\simeq}
    L_{B/A}.
\]

The transitivity triangle now gives
\[
    L_{B/C}
    \simeq
    0.
\]
Together with the isomorphism
\[
    \pi_0(C)
    \xrightarrow{\simeq}
    \pi_0(B),
\]
cotangent rigidity
\Ref{cor:cotangent-rigidity}
implies that
\[
    C
    \xrightarrow{\simeq}
    B.
\]
Thus, in a Zariski neighbourhood of the selected point, the original
closed immersion is the derived zero locus of
\[
    f_1,\ldots,f_n.
\]
This proves
\[
    \textup{(i)}
    \Longrightarrow
    \textup{(ii)}
\]
and completes the proof.
\end{proof}

\section{The deformation-theoretic output}
\label{sec:spectral-deformation-output}

The constructions of this chapter equip the spectral schemes of
\Ref{chap:spectral-algebraic-geometry} with relative and absolute
cotangent complexes, structured square-zero extensions, multiplicative
Postnikov $k$-invariants, affine and global obstruction theories, and the
local geometry of quasi-smooth closed immersions. All global objects lie
in the quasi-coherent categories and use the pullback, base-change,
relative-spectrum, closed-immersion, and linear-space constructions of
the preceding chapter. Smooth, \`etale, and Nisnevich geometry are
developed in \Ref{chap:smooth-spectral-motivic-site}.
